\documentclass[10pt, letterpaper]{article}

% Inhaltsverzeichnis für Pakettypen (nur für Übersicht im Header, wird nicht im Dokument angezeigt)
% 1. Seitenlayout und Ränder
% 2. Sprache und Zeichensatz
% 3. Mathematik und Theorem-Umgebungen
% 4. Eigene Makros
% 5. Diagramme und Grafiken
% 6. Tabellen und Aufzählungen
% 7. Inhaltsverzeichnis
% 8. Abschnittsüberschriften
% 9. Abstrakt-Umgebung
% 10. Todos/Notizen
% 11. Rahmen/Box-Umgebungen
% 12. Python-Integration
% 13. Literaturverwaltung
% 14. Hyperlinks
% 15. Absatzeinstellungen
% 16. Umgebungen
% 17  Graphik
% 18  Extra
% 00. Titel und Autor

\usepackage{slashed}

% --- 1. Seitenlayout und Ränder ---
\usepackage[margin=3cm]{geometry}

% --- 2. Sprache und Zeichensatz ---
\usepackage[english]{babel}
\usepackage[T1]{fontenc}
\usepackage[utf8]{inputenc}

% --- 3. Mathematik und Theorem-Umgebungen ---
\usepackage{amsmath, amssymb, amsthm}
\usepackage{mathrsfs}
\DeclareMathOperator{\WF}{WF}

% --- 4. Eigene Makros ---
\usepackage{xcolor}
\newcommand{\SKP}{\langle\cdot,\cdot\rangle}
\newcommand{\R}{\mathbb{R}}
\newcommand{\N}{\mathbb{N}}
\newcommand{\Q}{\mathbb{Q}}
\newcommand{\Z}{\mathbb{Z}}
\newcommand{\C}{\mathbb{C}}
\newcommand{\entwurf}[1]{\textcolor{red}{#1}}
\newcommand{\GL}{\mathrm{GL}}
\newcommand{\SL}{\mathrm{SL}}
\newcommand{\SO}{\mathrm{SO}}
\newcommand{\SU}{\mathrm{SU}}
\newcommand{\Sp}{\mathrm{Sp}}
\newcommand{\Ogroup}{\mathrm{O}}
\newcommand{\U}{\mathrm{U}}

% --- 5. Diagramme und Grafiken ---
\usepackage{graphicx}
\graphicspath{{../../Library/Images/}}
\usepackage[export]{adjustbox}
\usepackage{tikz}
\usetikzlibrary{decorations.pathreplacing, arrows.meta, positioning}
\usepackage{tikz-cd}

% --- 6. Tabellen und Aufzählungen ---
\usepackage{enumitem}
\setlist[itemize]{left=0.5cm}

\newenvironment{romanenum}[1][]
  {%
    \ifx&#1&
    \else
      \textbf{#1}\quad
    \fi
    \begin{enumerate}[label=\roman*)]
  }
  {%
    \end{enumerate}%
  }

% --- 7. Inhaltsverzeichnis ---
\usepackage{tocloft}
\renewcommand{\cftsecfont}{\footnotesize}
\renewcommand{\cftsubsecfont}{\footnotesize}
\renewcommand{\cftsubsubsecfont}{\footnotesize}
\renewcommand{\cftsecpagefont}{\footnotesize}
\renewcommand{\cftsubsecpagefont}{\footnotesize}
\renewcommand{\cftsubsubsecpagefont}{\footnotesize}
\usepackage{etoc}

% --- 8. Abschnittsüberschriften ---
\usepackage{titlesec}
\titleformat{\section}{\normalfont\large\bfseries}{\thesection}{1em}{}
\titleformat{\subsection}{\normalfont\normalsize\bfseries}{\thesubsection}{0.5em}{}
\titleformat{\subsubsection}{\normalfont\normalsize\bfseries}{\thesubsubsection}{0.5em}{}
\setcounter{secnumdepth}{4}

% --- 9. Abstrakt-Umgebung ---
\usepackage{changepage}
\renewenvironment{abstract}
  {
    \begin{adjustwidth}{1.5cm}{1.5cm}
    \small
    \textsc{Abstract. –}%
  }
  {
    \end{adjustwidth}
  }

% --- 10. Todos/Notizen ---
\usepackage{todonotes}
% Hellblau definieren
\definecolor{hellblau}{rgb}{0.3,0.6,1}
\newenvironment{note}{\par\begingroup\color{hellblau}\itshape}{\par\endgroup}

% --- 11. Rahmen/Box-Umgebungen ---
\usepackage{mdframed}
\usepackage{tcolorbox}
\colorlet{shadecolor}{gray!25}

\newenvironment{customTheorem}
  {\vspace{10pt}%
   \begin{mdframed}[
     backgroundcolor=gray!20,
     linewidth=0pt,
     innertopmargin=10pt,
     innerbottommargin=10pt,
     skipabove=\dimexpr\topsep+\ht\strutbox\relax,
     skipbelow=\topsep,
   ]}
  {\end{mdframed}
   \vspace{10pt}%
  }

% --- 12. Python-Integration ---
% (Deaktiviert in dieser Version, aktiviere bei Bedarf)
% \usepackage{pythontex}
% \usepackage[makestderr]{pythontex}

% --- 13. Literaturverwaltung ---
\usepackage{csquotes}
\usepackage[backend=biber, style=alphabetic, citestyle=alphabetic]{biblatex}
\addbibresource{bibliography.bib}

% --- 14. Hyperlinks ---
\usepackage{hyperref}
\hypersetup{
  colorlinks   = true,
  urlcolor     = blue,
  linkcolor    = blue,
  citecolor    = blue,
  frenchlinks  = true
}

% --- 15. Absatzeinstellungen ---
\usepackage[parfill]{parskip}
\sloppy

% --- 16. Umgebungen ---
\usepackage{thmtools}

\newcommand{\CustomHeading}[3]{%
  \par\medskip\noindent%
  \textbf{#1 #2} \textnormal{(#3)}.\enskip%
}

\newenvironment{DEF}[2]{\begin{unitbox}\CustomHeading{Definition}{#1}{#2}}{\end{unitbox}}
\newenvironment{PROP}[2]{\begin{unitbox}\CustomHeading{Proposition}{#1}{#2}}{\end{unitbox}}
\newenvironment{THEO}[2]{\begin{unitbox}\CustomHeading{Theorem}{#1}{#2}}{\end{unitbox}}
\newenvironment{LEM}[2]{\begin{unitbox}\CustomHeading{Lemma}{#1}{#2}}{\end{unitbox}}
\newenvironment{KORO}[2]{\begin{unitbox}\CustomHeading{Corollar}{#1}{#2}}{\end{unitbox}}
\newenvironment{REM}[2]{\begin{unitbox}\CustomHeading{Remark}{#1}{#2}}{\end{unitbox}}
\newenvironment{EXA}[2]{\begin{unitbox}\CustomHeading{Example}{#1}{#2}}{\end{unitbox}}
\newenvironment{STUD}[2]{\begin{unitbox}\CustomHeading{Study}{#1}{#2}}{\end{unitbox}}
\newenvironment{CONC}[2]{\begin{unitbox}\CustomHeading{Concept}{#1}{#2}}{\end{unitbox}}
\newenvironment{OTH}[2]{\begin{unitbox}\CustomHeading{Other}{#1}{#2}}{\end{unitbox}}
\newenvironment{EXE}[2]{\begin{unitbox}\CustomHeading{Exercise}{#1}{#2}}{\end{unitbox}}
\newenvironment{MOT}[2]{\begin{unitbox}\CustomHeading{Motivation}{#1}{#2}}{\end{unitbox}}
\newenvironment{PROOF}[2]{\begin{unitbox}\CustomHeading{Proof}{#1}{#2}}{\end{unitbox}}

% --- Unit Umgebung für Source-Inhalte ---
\usepackage{mdframed}
\newmdenv[
  linewidth=1pt,
  topline=false,
  bottomline=false,
  rightline=false,
  leftmargin=0cm,
  rightmargin=0cm,
  skipabove=10pt,
  skipbelow=10pt,
  innertopmargin=0.5\baselineskip,
  innerbottommargin=0.5\baselineskip,
  backgroundcolor=gray!10,
  linecolor=gray
]{unitbox}

\newenvironment{unit}[1]
  {\begin{unitbox}\textbf{Unit #1}\par\smallskip}
  {\end{unitbox}}

% --- 17. ... ---

% --- 18. Extras ---
\usepackage{stmaryrd}
\usepackage{bbold}  % falls du \mathbb{1} nutzen willst

% --- 00. Titel und Autor ---
\title{Mein Titel}
\author{Tim Jaschik}
\date{\today}

\begin{document}

\maketitle
\rule{\textwidth}{0.5pt}
\begin{abstract}
Kurze Beschreibung …
\end{abstract}
\rule{\textwidth}{0.5pt}
\vspace{0.5cm}

\tableofcontents

\pagebreak

\section*{Matrixgruppen}



Sei $\mathbb{K}$ ein Körper, typischerweise $\mathbb{K}=\mathbb{R}$ oder $\mathbb{K}=\mathbb{C} . M_{m, n}(\mathbb{K})$ ist die Menge der $m \times n$-Matrizen mit Einträgen in $\mathbb{K} . M_{n}(\mathbb{K})=M_{n, n}(\mathbb{K})$ ist die Menge der quadratischen Matrizen. Weiterhin ist $\mathbb{K}^{n}=M_{n, 1}(\mathbb{K})$. $M_{m, n}(\mathbb{K})$ ist ein Vektorraum der Dimension $m n$, eine Basis ist gegeben durch die Matrizen $E^{r, s}, 1 \leq r \leq m, 1 \leq s \leq n$, welche an der Stelle $(r, s)$ eine 1 stehen haben und sonst Nullen. Der Raum $M_{n}(\mathbb{K})$ ist abgeschlossen unter Matrixmultiplikation und bildet eine Algebra.

Die Determinantenfunktion det : $M_{n}(\mathbb{K}) \rightarrow \mathbb{K}$ hat folgende Eigenschaften:\\
(1) $\operatorname{det}(A B)=\operatorname{det} A \cdot \operatorname{det} B$.\\
(2) $\operatorname{det} I_{n}=1$, wobei $I_{n}$ die Einheitsmatrix ist.\\
(3) $A \in M_{n}(\mathbb{K})$ ist invertierbar genau dann, wenn $\operatorname{det} A \neq 0$.


\begin{DEF}{LIE-B26-01-01}{Matrizenräume und Determinante}
Sei $\mathbb{K}$ ein Körper, typischerweise $\mathbb{K}=\mathbb{R}$ oder $\mathbb{K}=\mathbb{C}$.  

\textbf{(a) Matrizenräume.}  
\begin{itemize}
  \item $M_{m,n}(\mathbb{K})$ bezeichnet den $\mathbb{K}$-Vektorraum aller $m\times n$-Matrizen mit Einträgen in $\mathbb{K}$.
  \item Für $n\in\mathbb{N}$ gilt $M_{n}(\mathbb{K}) := M_{n,n}(\mathbb{K})$ (Raum der quadratischen Matrizen).
  \item Der Spaltenvektorraum $\mathbb{K}^{n}$ wird identifiziert mit $M_{n,1}(\mathbb{K})$.
\end{itemize}
Der Raum $M_{m,n}(\mathbb{K})$ hat Dimension $mn$. Eine kanonische Basis ist gegeben durch die Matrizen
\[
E^{r,s},\quad 1\le r\le m,\ 1\le s\le n,
\]
wobei $E^{r,s}$ an der Stelle $(r,s)$ den Eintrag $1$ und sonst $0$ besitzt:
\[
(E^{r,s})_{ij} = \delta_{ir}\delta_{js}.
\]

\textbf{(b) Algebraeigenschaft.}  
Der Raum $M_{n}(\mathbb{K})$ ist unter der Matrixmultiplikation abgeschlossen und bildet mit dieser Operation eine \textbf{$\mathbb{K}$-Algebra}.

\textbf{(c) Determinantenfunktion.}  
Die Determinantenabbildung
\[
\det : M_{n}(\mathbb{K}) \longrightarrow \mathbb{K}
\]
erfüllt die folgenden fundamentalen Eigenschaften:
\begin{enumerate}
  \item \textbf{Multiplikativität:} $\det(AB) = (\det A)\cdot(\det B)$ für alle $A,B\in M_{n}(\mathbb{K})$.
  \item \textbf{Normierung:} $\det(I_n)=1$, wobei $I_n$ die $n\times n$-Einheitsmatrix ist.
  \item \textbf{Kriterium der Invertierbarkeit:}  
  Eine Matrix $A\in M_{n}(\mathbb{K})$ ist genau dann invertierbar, wenn $\det(A)\ne 0$ gilt.
\end{enumerate}
\end{DEF}


Definition 1. 

(1) $\mathrm{GL}_{n}(\mathbb{K})=\left\{A \in M_{n}(\mathbb{K}): \operatorname{det} A \neq 0\right\}$ ist die allgemeine lineare Gruppe.\\
(2) $\mathrm{SL}_{n}(\mathbb{K})=\left\{A \in M_{n}(\mathbb{K}): \operatorname{det} A=1\right\}$ ist die spezielle lineare Gruppe.\\
Offensichtlich sind $\mathrm{GL}_{n}(\mathbb{K})$ und $\mathrm{SL}_{n}(\mathbb{K})$ Gruppen. $\mathrm{SL}_{n}(\mathbb{K})$ ist eine Untergruppe von $\mathrm{GL}_{n}(\mathbb{K})$.

Definition 2. Sei $\mathbb{K}=\mathbb{R}$ oder $\mathbb{C}$.\\
(1) Ist $x=\left(x_{1}, \ldots, x_{n}\right)^{T} \in \mathbb{K}^{n}$, so setzen wir

$$
|x|=\sqrt{\left|x_{1}\right|^{2}+\ldots+\left|x_{n}\right|^{2}}
$$

(2) Ist $A \in M_{n}(\mathbb{K})$, so setzen wir

$$
\|A\|:=\sup \left\{|A x|: x \in \mathbb{K}^{n},|x|=1\right\}
$$

(3) $\rho(A, B):=\|A-B\|$ definiert eine Metrik und somit auch eine Topologie auf $\mathrm{GL}_{n}(\mathbb{K})$.



\begin{EXA}{LIE-B26-01-02}{Allgemeine und spezielle lineare Gruppe}
\textbf{(1) Allgemeine lineare Gruppe.}  
Für einen Körper $\mathbb{K}$ definieren wir
\[
\mathrm{GL}_{n}(\mathbb{K}) := \{\,A \in M_{n}(\mathbb{K}) \mid \det(A) \neq 0\,\}.
\]
$\mathrm{GL}_{n}(\mathbb{K})$ ist die \textbf{allgemeine lineare Gruppe} der invertierbaren $n\times n$-Matrizen über $\mathbb{K}$.

\textbf{(2) Spezielle lineare Gruppe.}  
\[
\mathrm{SL}_{n}(\mathbb{K}) := \{\,A \in M_{n}(\mathbb{K}) \mid \det(A) = 1\,\}.
\]
$\mathrm{SL}_{n}(\mathbb{K})$ ist die \textbf{spezielle lineare Gruppe} und eine Untergruppe von $\mathrm{GL}_{n}(\mathbb{K})$.
\end{EXA}



\begin{DEF}{LIE-B26-01-03}{Norm und Topologie auf Matrixgruppen}
\textbf{(1) Normen auf $\mathbb{K}^{n}$ und $M_{n}(\mathbb{K})$.}\\
Sei $\mathbb{K}=\mathbb{R}$ oder $\mathbb{C}$.
\begin{enumerate}
  \item Für $x=(x_{1},\ldots,x_{n})^{T}\in \mathbb{K}^{n}$ definieren wir die euklidische Norm durch
  \[
  |x| := \sqrt{|x_{1}|^{2}+\dots+|x_{n}|^{2}}.
  \]
  \item Für $A\in M_{n}(\mathbb{K})$ setzen wir die induzierte Operatornorm
  \[
  \|A\| := \sup\{\,|A x| : x\in \mathbb{K}^{n},\ |x|=1\,\}.
  \]
  Sie misst den größten Streckungsfaktor, den $A$ auf die Einheitskugel ausübt.
\end{enumerate}

\textbf{(2) Metrische Struktur.}  
Die Abbildung
\[
\rho(A,B):=\|A-B\|,\quad A,B\in M_{n}(\mathbb{K}),
\]
definiert eine Metrik auf $M_{n}(\mathbb{K})$ und induziert insbesondere eine \textbf{Topologie auf} $\mathrm{GL}_{n}(\mathbb{K})$ (als offene Teilmenge von $M_{n}(\mathbb{K})$).  
Damit wird $\mathrm{GL}_{n}(\mathbb{K})$ ein topologischer Gruppenraum.

Bemerkung: da $M_{n}(\mathbb{K})$ ein endlichdimensionaler Vektorraum ist, sind alle Normen äquivalent. Wir wählen diese spezielle Norm wegen ihrer im Lemma genannten Eigenschaften.
\end{DEF}



Lemma 3. 


\begin{LEM}{LIE-B26-01-04}{Elementare Eigenschaften der Operatornorm}
Es gilt $\|A \cdot B\| \leq\|A\| \cdot\|B\|$ und $\left\|I_{n}\right\|=1$.
\end{LEM}

Beweis.

\begin{PROOF}{LIE-B26-01-05}{P: Elementare Eigenschaften der Operatornorm}
$$
\begin{aligned}
\|A \cdot B\| & =\sup \{|A \cdot B \cdot x|:|x|=1\} \\
& =\sup \left\{\frac{|A \cdot B \cdot x|}{|B \cdot x|} \cdot|B \cdot x|:|x|=1\right\} \\
& =\sup \{[|A \cdot \underbrace{\frac{B x}{|B x|}}_{|\bullet|=1}| \cdot|B x|:|x|=1\} \\
& \leq \sup \{|A \cdot y|:|y|=1\} \cdot \sup \{[B x|:|x| \leq 1\} \\
& =\|A\| \cdot\|B\| .
\end{aligned}
$$
\end{PROOF}



Lemma 4. 


\begin{LEM}{LIE-B26-01-06}{Koordinantenfunktion, Determinante und Spur sind stetig auf $M_{n}(\mathbb{K})$}
Die Koordinatenfunktionen $M_{n}(\mathbb{K}) \rightarrow \mathbb{K}, A \mapsto A_{r s}, 1 \leq r, s \leq n$ sind stetig. Insbesondere sind $\operatorname{det}: M_{n}(\mathbb{K}) \rightarrow \mathbb{K}$ und $\operatorname{tr}: M_{n}(\mathbb{K}) \rightarrow \mathbb{K}, A \mapsto \sum_{i=1}^{n} A_{i i}$ stetig.
\end{LEM}

Beweis. 

\begin{PROOF}{LIE-B26-01-07}{P: Koordinantenfunktion, Determinante und Spur sind stetig auf $M_{n}(\mathbb{K})$}
Sei $e_{i}$ der $i$-te Standardbasisvektor. Dann ist
$$
\left|A_{r s}\right| \leq \sqrt{\sum_{i=1}^{n}\left|A_{i s}\right|^{2}}=\left|\sum_{i=1}^{n} A_{i s} e_{i}\right|=\left|A e_{s}\right| \leq\|A\|
$$
also folgt $\left|A_{r s}^{\prime}-A_{r s}\right| \leq\left\|A^{\prime}-A\right\|$.
\end{PROOF}



Definition 5. 

Sei $G$ eine Gruppe, versehen mit einer Topologie. Dann heisst $G$ eine topologische Gruppe, wenn die Multiplikation $m$ : $G \times G \rightarrow G,(g, h) \mapsto g h$ und die Inversenbildung $i: G \rightarrow G, g \mapsto g^{-1}$ stetig sind.


\begin{DEF}{LIE-B26-01-08}{Topologische Gruppe}
Sei $G$ eine Gruppe, versehen mit einer Topologie.  
Dann heißt $G$ eine \textbf{topologische Gruppe}, wenn die folgenden beiden Abbildungen stetig sind:
\begin{align*}
m &: G \times G \longrightarrow G, \quad (g,h) \longmapsto g h 
&& \text{(\textit{Multiplikation})},\\[3pt]
i &: G \longrightarrow G, \quad g \longmapsto g^{-1} 
&& \text{(\textit{Inversenbildung})}.
\end{align*}
Das heißt, sowohl die Verknüpfung zweier Gruppen­elemente als auch das Bilden des Inversen sind stetige Operationen bezüglich der gegebenen Topologie auf $G$.

\textbf{Bemerkung.}  
Eine topologische Gruppe ist damit gleichzeitig ein \textit{Gruppenobjekt} in der Kategorie der topologischen Räume:  
Die algebraische und die topologische Struktur sind kompatibel.
\end{DEF}


Lemma 6. 


\begin{EXA}{LIE-B26-01-09}{$\mathrm{GL}_{n}(\mathbb{K}), \mathrm{SL}_{n}(\mathbb{K})$ sind topologische Gruppen}
$\mathrm{GL}_{n}(\mathbb{K})$ und $\mathrm{SL}_{n}(\mathbb{K})$ sind topologische Gruppen.

Beweis. Stetitgkeit von $m$ ist klar. Stetigkeit von $i$ folgt aus der Tatsache, dass jeder Eintrag von $A^{-1}$ von der Gestalt
$$
(\operatorname{det} A)^{-1} \cdot \text { Polynom in den Einträgen von } A
$$
ist.
\end{EXA}


Definition 7. 

(1) Eine abgeschlossene Untergruppe von $\mathrm{GL}_{n}(\mathbb{K})$ heisst Matrixgruppe.\\
(2) Ist $G$ eine Matrixgruppe und $H \subset G$ eine abgeschlossene Untergruppe, so heisst $H$ Matrixuntergruppe von $G$.

Bemerkung: Ist $K$ Matrixuntergruppe von $G$ und $H$ Matrixuntergruppe von $K$, so ist $H$ auch Matrixuntergruppe von $G$.


\begin{DEF}{LIE-B26-01-10}{Matrixgruppen und Matrixuntergruppen}
\textbf{(1) Matrixgruppe.}  
Eine \textbf{Matrixgruppe} ist eine \emph{abgeschlossene Untergruppe} von $\mathrm{GL}_{n}(\mathbb{K})$, d.\,h.
\[
G \subseteq \mathrm{GL}_{n}(\mathbb{K}), \qquad G \text{ ist eine Gruppe und abgeschlossen in } \mathrm{GL}_{n}(\mathbb{K})
\]
(bezüglich der von der Operatornorm induzierten Topologie).

\textbf{(2) Matrixuntergruppe.}  
Sei $G$ eine Matrixgruppe.  
Eine \textbf{Matrixuntergruppe} von $G$ ist eine abgeschlossene Untergruppe $H \subseteq G$.

Ist $K$ eine Matrixuntergruppe von $G$ und $H$ eine Matrixuntergruppe von $K$,  
so ist $H$ ebenfalls eine Matrixuntergruppe von $G$.  
Mit anderen Worten: die Eigenschaft „Matrixuntergruppe“ ist \emph{transitiv}.
\end{DEF}




\section*{Beispiele:}


(1) $\mathrm{GL}_{n}(\mathbb{K})$ ist eine Matrixgruppe. Achtung: als Teilmenge von $M_{n}(\mathbb{K})$ ist $\mathrm{GL}_{n}(\mathbb{K})$ nicht abgeschlossen, sondern offen. Als Teilmenge von sich selbst ist es sowohl offen als auch abgeschlossen.\\
(2) $\mathrm{SL}_{n}(\mathbb{K})$ ist eine Matrixgruppe.\\
(3) Unter der Abbildung $A \mapsto\left(\begin{array}{cc}A & 0 \\ 0 & 1\end{array}\right)$ wird $\mathrm{GL}_{n}(\mathbb{K})$ zu einer Matrixuntergruppe von $\mathrm{GL}_{n+1}(\mathbb{K})$ und $\mathrm{SL}_{n}(\mathbb{K})$ zu einer Matrixuntergruppe von $\mathrm{SL}_{n+1}(\mathbb{K})$.\\
(4) $\mathrm{UT}_{n}(\mathbb{K})$ die Gruppe der oberen triangulären Matrizen.\\
(5) $\operatorname{SUT}_{n}(\mathbb{K})$ die Gruppe der unipotenten oberen triangulären Matrizen, d.h. derjenigen, bei denen auf der Diagonalen nur Einsen stehen.\\
(6) $\theta: \mathbb{K} \rightarrow \operatorname{SUT}_{2}(\mathbb{K}), t \mapsto\left(\begin{array}{cc}1 & t \\ 0 & 1\end{array}\right)$ ist ein Homomorphismus und ein Gruppenisomorphismus. Also ist ( $\mathbb{K},+$ ) eine Matrixgruppe.\\
(7)

$$
\operatorname{Aff}_{n}(\mathbb{K})=\left\{\left(\begin{array}{cc}
A & t \\
0 & 1
\end{array}\right): A \in \mathrm{GL}_{n}(\mathbb{K}), t \in \mathbb{K}^{n}\right\}
$$

ist eine Matrixgruppe, die affine Gruppe genannt wird. Sie operiert auf $\mathbb{K}^{n}$ durch affine Transformationen $(A, t) x:=A x+ t$.\\



\begin{EXA}{LIE-B26-01-11}{Allgemeine lineare Gruppe als Matrixgruppe}
$\mathrm{GL}_{n}(\mathbb{K})$ ist eine \textbf{Matrixgruppe}.  
Sie ist als Untergruppe von $(M_{n}(\mathbb{K}),\cdot)$ abgeschlossen unter Matrixmultiplikation und Inversenbildung.  
Als Teilmenge von $M_{n}(\mathbb{K})$ ist $\mathrm{GL}_{n}(\mathbb{K})$ \emph{offen}, da die Determinantenabbildung $\det : M_{n}(\mathbb{K}) \to \mathbb{K}$ stetig ist und $\mathrm{GL}_{n}(\mathbb{K}) = \det^{-1}(\mathbb{K}\setminus\{0\})$.  
Als Teilmenge von sich selbst ist sie jedoch sowohl offen als auch abgeschlossen.
\end{EXA}



\begin{EXA}{LIE-B26-01-12}{Spezielle lineare Gruppe}
$\mathrm{SL}_{n}(\mathbb{K}) := \{A \in \mathrm{GL}_{n}(\mathbb{K}) : \det A = 1\}$  
ist eine \textbf{Matrixgruppe}.  
Sie ist eine abgeschlossene Untergruppe von $\mathrm{GL}_{n}(\mathbb{K})$, da $\mathrm{SL}_{n}(\mathbb{K}) = \det^{-1}(\{1\})$ und $\det$ stetig ist.
\end{EXA}



\begin{EXA}{LIE-B26-01-13}{Einbettung von $\GL,\SL$ in höhere Dimensionen}
Unter der Abbildung
\[
A \longmapsto 
\begin{pmatrix}
A & 0 \\[3pt]
0 & 1
\end{pmatrix}
\]
wird $\mathrm{GL}_{n}(\mathbb{K})$ zu einer \textbf{Matrixuntergruppe} von $\mathrm{GL}_{n+1}(\mathbb{K})$  
und $\mathrm{SL}_{n}(\mathbb{K})$ zu einer Matrixuntergruppe von $\mathrm{SL}_{n+1}(\mathbb{K})$.
\end{EXA}



\begin{EXA}{LIE-B26-01-14}{$\operatorname{UT}$ Gruppe der oberen Dreiecksmatrizen}
\[
\mathrm{UT}_{n}(\mathbb{K}) := 
\left\{
A = (a_{ij}) \in \mathrm{GL}_{n}(\mathbb{K}) : a_{ij} = 0 \text{ für } i > j
\right\}
\]
ist eine \textbf{Matrixgruppe}.  
Sie ist abgeschlossen in $\mathrm{GL}_{n}(\mathbb{K})$, da die Bedingung $a_{ij}=0$ für $i>j$ durch stetige Gleichungen beschrieben wird.
\end{EXA}



\begin{EXA}{LIE-B26-01-15}{$\operatorname{SUT}$ unipotente obere Dreiecksmatrixgruppe}
\[
\mathrm{SUT}_{n}(\mathbb{K}) := 
\left\{
A = (a_{ij}) \in \mathrm{UT}_{n}(\mathbb{K}) : a_{ii} = 1 \text{ für alle } i
\right\}
\]
heißt \textbf{unipotente obere Dreiecksmatrixgruppe}.  
Sie ist ebenfalls abgeschlossen, da sie durch die zusätzlichen Gleichungen $a_{ii}=1$ charakterisiert ist.
\end{EXA}



\begin{EXA}{LIE-B26-01-16}{Additive Gruppe als Matrixgruppe}
Die Abbildung
\[
\theta: \mathbb{K} \longrightarrow \mathrm{SUT}_{2}(\mathbb{K}), 
\quad t \longmapsto 
\begin{pmatrix}
1 & t \\[3pt]
0 & 1
\end{pmatrix}
\]
ist ein \textbf{Gruppenisomorphismus}.  
Damit ist $(\mathbb{K},+)$ als additive Gruppe isomorph zu einer Matrixgruppe.
\end{EXA}



\begin{EXA}{LIE-B26-01-17}{Affine Gruppe}
\[
\mathrm{Aff}_{n}(\mathbb{K}) :=
\left\{
\begin{pmatrix}
A & t \\[3pt]
0 & 1
\end{pmatrix}
: A \in \mathrm{GL}_{n}(\mathbb{K}),\ t \in \mathbb{K}^{n}
\right\}
\]
ist eine \textbf{Matrixgruppe}, genannt die \emph{affine Gruppe}.  
Sie operiert auf $\mathbb{K}^{n}$ durch affine Transformationen:
\[
(A,t)\cdot x := A x + t.
\]
\end{EXA}





(8) $\mathrm{O}(n):=\left\{A \in \mathrm{GL}_{n}(\mathbb{R}): A^{t} A=I_{n}\right\} \subset \mathrm{GL}_{n}(\mathbb{R})$ die orthogonale Gruppe.\\
(9) $\mathrm{SO}(n):=\{A \in \mathrm{O}(n): \operatorname{det} A=1\}$ die spezielle orthogonale Gruppe.\\
(10) $\operatorname{Isom}_{n}(\mathbb{R})=\left\{\left(\begin{array}{cc}A & t \\ 0 & 1\end{array}\right): A \in \mathrm{O}_{n}(\mathbb{R}), t \in \mathbb{R}^{n}\right\}$ die Isometriegruppe von $\mathbb{R}^{n}$.\\
(11)

$$
\text { Lor }:=\left\{A \in \mathrm{GL}_{4}(\mathbb{R}): A t Q A=Q\right\}
$$

wobei $Q:=\left(\begin{array}{cccc}1 & 0 & 0 & 0 \\ 0 & 1 & 0 & 0 \\ 0 & 0 & 1 & 0 \\ 0 & 0 & 0 & -1\end{array}\right)$ heisst Lorentzgruppe.\\
(12) Sei $J=\left(\begin{array}{cc}0 & 1 \\ -1 & 0\end{array}\right)$ und $J_{2 m}$ die Blockdiagonalmatrix mit $2 \times 2$-Blöcken $J$ auf der Diagonale. Die reelle symplektische Gruppe ist

$$
\operatorname{Symp}_{2 n}(\mathbb{R}):=\left\{A \in \mathrm{GL}_{n}(\mathbb{R}): A^{t} J_{2 m} A=J_{2 m}\right\}
$$

(13) $\mathrm{U}(n):=\left\{A \in \mathrm{GL}_{n}(\mathbb{C}): A^{*} A=I_{n}\right\} \subset \mathrm{GL}_{n}(\mathbb{C})$ die unitäre Gruppe.\\
(14) $\operatorname{SU}(n):=\left\{A \in \mathrm{U}_{n}(\mathbb{C}): \operatorname{det} A=1\right\}$ die spezielle unitäre Gruppe.\\



\begin{EXA}{LIE-B26-01-18}{Orthogonale Gruppe}
Die \textbf{orthogonale Gruppe} ist definiert durch
\[
\mathrm{O}(n)
:= \left\{ A \in \mathrm{GL}_{n}(\mathbb{R}) \; : \; A^{\mathsf{T}} A = I_{n} \right\}.
\]
Sie besteht aus allen reellen Matrizen, die das Standardskalarprodukt auf $\mathbb{R}^{n}$ erhalten.  
$\mathrm{O}(n)$ ist eine abgeschlossene Untergruppe von $\mathrm{GL}_{n}(\mathbb{R})$, da die Bedingung $A^{\mathsf{T}}A = I_{n}$ stetig ist.
\end{EXA}



\begin{EXA}{LIE-B26-01-19}{Spezielle orthogonale Gruppe}
Die \textbf{spezielle orthogonale Gruppe} ist
\[
\mathrm{SO}(n)
:= \left\{ A \in \mathrm{O}(n) \; : \; \det A = 1 \right\}.
\]
Sie besteht aus den orientierungserhaltenden orthogonalen Transformationen und ist eine abgeschlossene Untergruppe von $\mathrm{O}(n)$.
\end{EXA}



\begin{EXA}{LIE-B26-01-20}{Isometriegruppe des euklidischen Raumes}
Die \textbf{Isometriegruppe} des $\mathbb{R}^{n}$ ist
\[
\mathrm{Isom}_{n}(\mathbb{R})
:= \left\{
\begin{pmatrix}
A & t \\[3pt]
0 & 1
\end{pmatrix}
: A \in \mathrm{O}(n),\ t \in \mathbb{R}^{n}
\right\}.
\]
Diese Gruppe beschreibt alle euklidischen Bewegungen (Drehungen, Spiegelungen und Translationen).  
Sie operiert auf $\mathbb{R}^{n}$ durch $(A,t)\cdot x = A x + t$.
\end{EXA}



\begin{EXA}{LIE-B26-01-21}{Lorentzgruppe}
Die \textbf{Lorentzgruppe} ist definiert als
\[
\mathrm{Lor}
:= \left\{
A \in \mathrm{GL}_{4}(\mathbb{R}) \; : \; A^{\mathsf{T}} Q A = Q
\right\},
\qquad
Q :=
\begin{pmatrix}
1 & 0 & 0 & 0 \\
0 & 1 & 0 & 0 \\
0 & 0 & 1 & 0 \\
0 & 0 & 0 & -1
\end{pmatrix}.
\]
Sie besteht aus allen linearen Transformationen des Minkowski-Raums $\mathbb{R}^{1,3}$, die die Lorentz-Metrik erhalten.  
$\mathrm{Lor}$ ist eine Matrixgruppe, die in der Relativitätstheorie zentrale Bedeutung hat.
\end{EXA}



\begin{EXA}{LIE-B26-01-22}{Reelle symplektische Gruppe}
Sei $J = 
\begin{pmatrix}
0 & 1 \\[3pt]
-1 & 0
\end{pmatrix}$ 
und $J_{2m}$ die Blockdiagonalmatrix mit $m$ Blöcken $J$ auf der Diagonale.  
Dann heißt
\[
\mathrm{Symp}_{2m}(\mathbb{R})
:= \left\{ A \in \mathrm{GL}_{2m}(\mathbb{R}) \; : \; A^{\mathsf{T}} J_{2m} A = J_{2m} \right\}
\]
die \textbf{reelle symplektische Gruppe}.  
Sie besteht aus allen linearen Transformationen, die die Standard-Symplektikform auf $\mathbb{R}^{2m}$ erhalten.
\end{EXA}



\begin{EXA}{LIE-B26-01-23}{Unitäre Gruppe}
Die \textbf{unitäre Gruppe} ist definiert durch
\[
\mathrm{U}(n)
:= \left\{ A \in \mathrm{GL}_{n}(\mathbb{C}) \; : \; A^{*} A = I_{n} \right\},
\]
wobei $A^{*} = \overline{A}^{\mathsf{T}}$ die adjungierte Matrix ist.  
$\mathrm{U}(n)$ besteht aus allen linearen Transformationen, die das Standardhermitesche Skalarprodukt auf $\mathbb{C}^{n}$ erhalten.
\end{EXA}



\begin{EXA}{LIE-B26-01-24}{Spezielle unitäre Gruppe}
Die \textbf{spezielle unitäre Gruppe} ist
\[
\mathrm{SU}(n)
:= \left\{ A \in \mathrm{U}(n) \; : \; \det A = 1 \right\}.
\]
Sie besteht aus den orientierungserhaltenden unitären Transformationen.  
$\mathrm{SU}(n)$ ist eine abgeschlossene Untergruppe von $\mathrm{U}(n)$ und spielt eine fundamentale Rolle in der Differentialgeometrie und in der theoretischen Physik (z.\,B. als Eichgruppe in der Quantenfeldtheorie).
\end{EXA}




Bemerkung: komplexe Matrixgruppen wie $\mathrm{U}(n)$ kann man auch als reelle auffassen. Definiere dazu $\rho_{n}^{\prime}: M_{n}(\mathbb{C}) \rightarrow M_{2 n}(\mathbb{R})$ durch $X+i Y \mapsto \left(\begin{array}{cc}X & -Y \\ Y & X\end{array}\right)$. Dann ist $\rho_{n}^{\prime} M_{n}(\mathbb{C})=\left\{A \in M_{2 n}(\mathbb{R}): A J_{2 n}^{\prime}=J_{2 n}^{\prime} A\right\}$, wobei $J_{2 n}^{\prime}=\left(\begin{array}{cc}0 & I_{n} \\ -I_{n} & 0\end{array}\right)$. Ist nämlich $A=\left(\begin{array}{cc}A_{11} & A_{12} \\ A_{21} & A_{22}\end{array}\right)$, so ist $A J_{2 n}^{\prime}=J_{2 n}^{\prime} A$ genau dann, wenn $\left(\begin{array}{cc}-A_{12} & A_{11} \\ -A_{22} & A_{21}\end{array}\right)=\left(\begin{array}{cc}A_{21} & A_{22} \\ -A_{11} & -A_{12}\end{array}\right)$



\begin{CONC}{LIE-B26-01-25}{Reelle Realisierung komplexer Matrixgruppen}
Komplexe Matrixgruppen wie etwa $\mathrm{U}(n)$ können auch als reelle Matrixgruppen aufgefasst werden.  
Dazu definiert man die Abbildung
\[
\rho_n' : M_n(\mathbb{C}) \longrightarrow M_{2n}(\mathbb{R}), \qquad
X + iY \longmapsto
\begin{pmatrix}
X & -Y \\[3pt]
Y & X
\end{pmatrix},
\]
wobei $X, Y \in M_n(\mathbb{R})$ sind.

\textbf{Eigenschaften:}
\begin{enumerate}
\item $\rho_n'$ ist ein reeller Algebrenhomomorphismus, d.\,h.
\[
\rho_n'(A + B) = \rho_n'(A) + \rho_n'(B), \qquad
\rho_n'(A B) = \rho_n'(A)\rho_n'(B).
\]
\item Das Bild von $\rho_n'$ ist gegeben durch
\[
\rho_n'(M_n(\mathbb{C})) =
\left\{
A \in M_{2n}(\mathbb{R}) \; : \; A J_{2n}' = J_{2n}' A
\right\},
\]
wobei
\[
J_{2n}' =
\begin{pmatrix}
0 & I_n \\[3pt]
- I_n & 0
\end{pmatrix}.
\]
\item Ist $A = 
\begin{pmatrix}
A_{11} & A_{12} \\[3pt]
A_{21} & A_{22}
\end{pmatrix}
\in M_{2n}(\mathbb{R})$, so gilt
\[
A J_{2n}' = J_{2n}' A
\quad \Longleftrightarrow \quad
\begin{pmatrix}
- A_{12} & A_{11} \\[3pt]
- A_{22} & A_{21}
\end{pmatrix}
=
\begin{pmatrix}
A_{21} & A_{22} \\[3pt]
- A_{11} & - A_{12}
\end{pmatrix}.
\]
Diese Bedingung charakterisiert also genau diejenigen reellen Matrizen, die aus komplexen Matrizen via $\rho_n'$ entstehen.
\end{enumerate}

\textbf{Interpretation:}  
Unter der Identifikation $\mathbb{C}^n \cong \mathbb{R}^{2n}$ mit komplexer Struktur $J_{2n}'$ wird eine komplexe Matrix $A \in M_n(\mathbb{C})$ durch $\rho_n'(A)$ als reelle lineare Abbildung beschrieben, die mit $J_{2n}'$ kommutiert.  
Somit kann jede komplexe Matrixgruppe, insbesondere $\mathrm{U}(n)$ oder $\mathrm{SU}(n)$, als reelle Matrixgruppe in $\mathrm{GL}_{2n}(\mathbb{R})$ aufgefasst werden.
\end{CONC}



Definition 8. 


Seien $G, H$ Matrixgruppen. Ein Gruppenhomomorphismus $\phi: G \rightarrow H$ ist ein stetiger Homomorphismus von Matrixgruppen, wenn $\phi$ stetig ist und das Bild $\phi(G)$ abgeschlossen in $H$ ist.


\begin{DEF}{LIE-B26-01-26}{Stetiger Homomorphismus von Matrixgruppen}
Seien $G$ und $H$ Matrixgruppen.  
Ein Gruppenhomomorphismus
\[
\phi : G \longrightarrow H
\]
heißt \textbf{stetiger Homomorphismus von Matrixgruppen}, wenn folgende Bedingungen erfüllt sind:
\begin{enumerate}
  \item $\phi$ ist \textbf{stetig} in der durch die Einbettung von $G, H \subseteq M_n(\mathbb{K})$ induzierten Topologie,
  \item das Bild $\phi(G)$ ist \textbf{abgeschlossen} in $H$.
\end{enumerate}
In diesem Fall schreibt man oft
\[
\phi : G \twoheadrightarrow \phi(G) \subseteq H
\]
und bezeichnet $\phi(G)$ als das \textbf{abgeschlossene Bild} von $\phi$ in $H$.  
Stetige Homomorphismen zwischen Matrixgruppen sind damit sowohl algebraisch (Gruppenhomomorphismen) als auch topologisch (stetige Abbildungen) verträglich mit der Struktur der Matrixgruppen.
\end{DEF}




Beispiel: Sei

$$
G:=\left\{\left(\begin{array}{cc}
1 & n \\
0 & 1
\end{array}\right): n \in \mathbb{Z}\right\}
$$

$\phi: G \rightarrow \mathrm{U}(1)=S^{1},\left(\begin{array}{cc}1 & n \\ 0 & 1\end{array}\right) \mapsto e^{2 \pi r n i}$, wobei $r$ eine feste irrationale Zahl ist. Dann ist $\phi(G)$ dicht in $\mathrm{U}(1)$, aber eine echte Teilmenge. Also ist $\phi$ kein stetiger Homomorphismus von Matrixgruppen.



\begin{EXA}{LIE-B26-01-27}{Nicht-abgeschlossenes Bild eines Homomorphismus}
Sei
\[
G := 
\left\{
\begin{pmatrix}
1 & n \\[3pt]
0 & 1
\end{pmatrix}
: n \in \mathbb{Z}
\right\}
\subset \mathrm{GL}_{2}(\mathbb{R})
\]
eine diskrete Matrixgruppe (isomorph zu $(\mathbb{Z}, +)$),
und definiere den Gruppenhomomorphismus
\[
\phi : G \longrightarrow \mathrm{U}(1) = S^{1}, 
\qquad
\begin{pmatrix}
1 & n \\[3pt]
0 & 1
\end{pmatrix}
\longmapsto e^{2 \pi i r n},
\]
wobei $r \in \mathbb{R}$ eine feste irrationale Zahl sei.
\smallskip

\textbf{Analyse:}
\begin{itemize}
  \item $\phi$ ist ein \emph{Gruppenhomomorphismus}, da 
  \[
  \phi\!\left(
  \begin{pmatrix}1 & n_1 \\ 0 & 1\end{pmatrix}
  \begin{pmatrix}1 & n_2 \\ 0 & 1\end{pmatrix}
  \right)
  = \phi\!\left(
  \begin{pmatrix}1 & n_1 + n_2 \\ 0 & 1\end{pmatrix}
  \right)
  = e^{2\pi i r (n_1+n_2)}
  = e^{2\pi i r n_1} e^{2\pi i r n_2}.
  \]
  \item Das Bild ist gegeben durch
  \[
  \phi(G) = \{ e^{2\pi i r n} : n \in \mathbb{Z} \}.
  \]
  Da $r$ irrational ist, ist die Menge 
  $\{ e^{2\pi i r n} : n \in \mathbb{Z} \}$
  \emph{dicht} in $S^1$, aber keine abgeschlossene Teilmenge.
\end{itemize}

\textbf{Folgerung:}
Das Bild $\phi(G)$ ist dicht, aber nicht abgeschlossen in $\mathrm{U}(1)$.  
Somit ist $\phi$ zwar ein stetiger Gruppenhomomorphismus, aber \emph{kein stetiger Homomorphismus von Matrixgruppen} im Sinne der Definition, da die Abgeschlossenheitsbedingung verletzt ist.
\end{EXA}


Definition 9. Ist $\phi: G \rightarrow H$ ein Homomorphismus von $M a$ trixgruppen und ist $\phi$ ein Homöomorphismus, dann heisst $\phi$ stetiger Isomorphismus von Matrixgruppen.



\begin{DEF}{LIE-B26-01-28}{Stetiger Isomorphismus von Matrixgruppen}
Seien $G$ und $H$ Matrixgruppen.  
Ein Gruppenhomomorphismus
\[
\phi : G \longrightarrow H
\]
heißt \textbf{stetiger Isomorphismus von Matrixgruppen},  
wenn gilt:
\begin{enumerate}
  \item $\phi$ ist ein \textbf{Homomorphismus von Matrixgruppen}, d.\,h. $\phi$ ist stetig und $\phi(G)$ ist abgeschlossen in $H$;
  \item $\phi$ ist ein \textbf{Homöomorphismus}, d.\,h. bijektiv, stetig und die Umkehrabbildung $\phi^{-1}$ ist ebenfalls stetig.
\end{enumerate}

In diesem Fall sind $G$ und $H$ nicht nur als Gruppen, sondern auch als topologische Räume (in der von $M_n(\mathbb{K})$ induzierten Topologie) isomorph.  
Man schreibt dann kurz
\[
G \cong_{\text{top}} H
\quad \text{oder} \quad
G \simeq H,
\]
und bezeichnet $\phi$ als \emph{stetigen Isomorphismus von Matrixgruppen}.
\end{DEF}



SATZ 10. 

Sei $\phi: G \rightarrow H$ ein stetiger Homomorphismus von $M a_{-}$ trixgruppen. Dann ist $\operatorname{ker} \phi$ eine Matrixuntergruppe von $G$ und die induzierte Abbildung $\bar{\phi}: G / \operatorname{ker} \phi \rightarrow \phi(G)$ ist ein Isomorphismus.

Beweis. Sei $\left(A_{n}\right)$ eine Folge in $\operatorname{ker} \phi$, die gegen $A \in G$ konvergiert. Da $\phi$ stetig ist, gilt $\phi(A)=\phi\left(\lim _{n} A_{n}\right)=\lim \phi\left(A_{n}\right)=1$. Die zweite Aussage folgt aus dem ersten Isomorphiesatz.


\begin{PROP}{LIE-B26-01-29}{Erster Isomorphiesatz für Matrixgruppen}
Sei $\phi: G \to H$ ein stetiger Homomorphismus von Matrixgruppen
(d.\,h. $\phi$ ist stetig und $\phi(G)$ ist abgeschlossen in $H$). Dann gilt:
\begin{enumerate}
  \item $\ker\phi$ ist eine \textbf{Matrixuntergruppe} von $G$ (insbesondere abgeschlossen und normal).
  \item Die induzierte Abbildung
  \[
    \bar\phi:\; G/\ker\phi \longrightarrow \phi(G),\qquad g\,\ker\phi \longmapsto \phi(g)
  \]
  ist ein \textbf{Isomorphismus von Matrixgruppen} (Gruppenisomorphismus und Homöomorphismus).
\end{enumerate}
\end{PROP}



\begin{PROOF}{LIE-B26-01-30}{P: Erster Isomorphiesatz für Matrixgruppen}
(1) \emph{$\ker\phi$ ist abgeschlossen und damit Matrixuntergruppe.}
Sei $(A_n)$ eine Folge in $\ker\phi$ mit $A_n \to A \in G$. Aus der Stetigkeit von $\phi$ folgt
\[
\phi(A)\;=\;\phi\!\big(\lim_n A_n\big)\;=\;\lim_n \phi(A_n)\;=\;\lim_n e\;=\;e,
\]
also $A\in\ker\phi$. Damit ist $\ker\phi$ abgeschlossen in $G$. Als Kern eines Homomorphismus ist $\ker\phi$ zudem normal; somit ist $\ker\phi$ eine Matrixuntergruppe von $G$.

\smallskip
(2) \emph{Algebraischer Isomorphismus.}
Nach dem Ersten Isomorphiesatz für Gruppen ist
\[
\bar\phi:\; G/\ker\phi \to \phi(G),\quad g\ker\phi \mapsto \phi(g)
\]
ein wohldefinierter bijektiver Gruppenhomomorphismus.

\smallskip
(3) \emph{Topologie und Stetigkeit.}
Da $\ker\phi$ abgeschlossen ist, ist der Quotient $G/\ker\phi$ ein Hausdorff’scher topologischer Gruppenraum; die kanonische Projektion
$\pi:G\to G/\ker\phi$ ist offen und stetig.
Die Abbildung $\bar\phi$ ist stetig, denn $\bar\phi\circ\pi=\phi$ und $\phi$ ist stetig, während $\pi$ surjektiv ist.

\smallskip
(4) \emph{$\bar\phi$ ist ein Homöomorphismus.}
Das Bild $\phi(G)$ ist nach Voraussetzung in $H$ abgeschlossen und damit eine Matrixgruppe. Die Umkehrabbildung
\[
\bar\phi^{-1}:\; \phi(G)\to G/\ker\phi,\qquad \bar\phi^{-1}(\phi(g))=g\ker\phi,
\]
ist stetig: Schreibe $\bar\phi^{-1} = \pi\circ s$, wobei $s:\phi(G)\to G$ eine (beliebige) Wahl einer Repräsentantenabbildung lokal verifiziert; allgemeiner folgt die Stetigkeit aus der universellen Eigenschaft des Quotienten: Ist $f:G\to \phi(G)$ stetig und konstant auf den Nebenklassen von $\ker\phi$ (hier $f=\phi$), so faktorisert $f$ eindeutig über eine stetige $\tilde f:G/\ker\phi\to \phi(G)$. Da $\bar\phi$ bijektiv ist, ist auch die inverse Faktorisierung stetig. Alternativ kann man verwenden, dass $\bar\phi$ eine stetige, offene Bijektion ist: Offenheit folgt, weil $\pi$ offen ist und $\phi=\bar\phi\circ\pi$ stetig-offen auf Nebenklassen konstant ist.

\smallskip
Damit ist $\bar\phi$ ein Gruppenisomorphismus und Homöomorphismus, also ein stetiger Isomorphismus von Matrixgruppen.
\end{PROOF}


Achtung: $G / \operatorname{ker} \phi$ ist im Allgemeinen keine Matrixgruppe. Das ist ein Grund, warum wir später Liegruppen studieren werden.

Definition 11. (1) Sei $G$ eine Gruppe und $X$ eine Menge. Eine Operation von $G$ auf $X$ ist eine Abbildung $\mu: G \times X \rightarrow X,(g, x) \mapsto g x$, so dass $(g h) x=g(h x), 1 x=x$ für alle $x \in G, g, h \in G$.\\
(2) Der Stabilisator im Punkt $x$ ist

$$
\operatorname{Stab}_{G}(x)=\{g \in G: g x=x\} \subset G .
$$

(3) Das Orbit von $x$ ist

$$
\operatorname{Orb}_{G}(x)=\{g x: g \in G\} \subset X
$$



\begin{DEF}{LIE-B26-07-01}{Gruppenoperation, Stabilisator und Orbit}
\begin{enumerate}
\item
Sei $G$ eine Gruppe und $X$ eine Menge.  
Eine \textbf{(Links-)Operation} von $G$ auf $X$ ist eine Abbildung
\[
\mu : G \times X \longrightarrow X,\qquad (g,x) \longmapsto g x,
\]
so dass für alle $g,h \in G$ und $x \in X$ gilt:
\begin{align*}
(\text{A1}) \quad &(g h) x = g (h x) &&\text{(Assoziativität der Wirkung)},\\
(\text{A2}) \quad &1 x = x &&\text{(Neutrales Element wirkt trivial).}
\end{align*}
In diesem Fall sagt man: $G$ \emph{operiert (von links)} auf $X$.  
Manchmal schreibt man $\mu(g,x)$ statt $g x$.

\item
Für ein festes $x \in X$ heißt
\[
\operatorname{Stab}_{G}(x)
:= \{\, g \in G : g x = x \,\}
\]
der \textbf{Stabilisator} (oder die \textbf{Isotropiegruppe}) von $x$.  
Er ist eine Untergruppe von $G$.

\item
Für ein festes $x \in X$ heißt
\[
\operatorname{Orb}_{G}(x)
:= \{\, g x : g \in G \,\}
\]
die \textbf{Bahn} (oder das \textbf{Orbit}) von $x$ unter der Wirkung von $G$.

\end{enumerate}
\end{DEF}





Satz 12. (1) $\operatorname{Stab}_{G}(x)$ ist eine Untergruppe von $G$.\\
(2) $y \in \operatorname{Orb}_{G}(x) \Longleftrightarrow \operatorname{Orb}_{G}(x)=\operatorname{Orb}_{G}(y)$.\\
(3) Für $x \in X$ gibt es eine Bijektion

$$
\phi: G / \operatorname{Stab}_{G}(x) \rightarrow \operatorname{Orb}_{G}(x),[g] \mapsto g x .
$$

(4) Ist $y=t x, t \in G$, so gilt

$$
\operatorname{Stab}_{G}(y)=t \operatorname{Stab}_{G}(x) t^{-1}
$$

Beweis. Nachrechnen.



\begin{PROP}{LIE-B26-07-02}{Elementare Eigenschaften zu Stabilisator, Bahn und Nebenklassen}
Sei $G$ eine Gruppe, die (von links) auf einer Menge $X$ operiert. Dann gilt für jedes $x\in X$:
\begin{enumerate}
\item $\operatorname{Stab}_G(x)$ ist eine Untergruppe von $G$.
\item Für $y\in X$ gilt: $y\in \operatorname{Orb}_G(x)$ genau dann, wenn $\operatorname{Orb}_G(y)=\operatorname{Orb}_G(x)$.
\item Die Abbildung
\[
\phi:\; G/\operatorname{Stab}_G(x)\longrightarrow \operatorname{Orb}_G(x),\qquad [g]\longmapsto g\cdot x,
\]
ist eine wohldefinierte Bijektion.
\item Ist $y=t\cdot x$ für ein $t\in G$, so gilt
\[
\operatorname{Stab}_G(y)=t\,\operatorname{Stab}_G(x)\,t^{-1}.
\]
\end{enumerate}
\end{PROP}



\begin{PROOF}{LIE-B26-07-03}{P: Elementare Eigenschaften zu Stabilisator, Bahn und Nebenklassen}
\textbf{(1) $\operatorname{Stab}_G(x)\le G$.}
(Neutral) $1\cdot x=x$, also $1\in\operatorname{Stab}_G(x)$. 
(Abgeschlossenheit) Sind $g,h\in\operatorname{Stab}_G(x)$, so $(gh)\cdot x=g\cdot(h\cdot x)=g\cdot x=x$, also $gh\in\operatorname{Stab}_G(x)$.
(Inverses) Ist $g\in\operatorname{Stab}_G(x)$, so $g\cdot x=x$. Wirkt man links mit $g^{-1}$, erhält man $x=g^{-1}\cdot x$, also $g^{-1}\in\operatorname{Stab}_G(x)$.

\smallskip
\textbf{(2) Gleiche Bahnen genau bei Erreichbarkeit.}
“Ist-Richtung”: $y\in\operatorname{Orb}_G(x)$, d.\,h. $y=g\cdot x$ für ein $g$. Für jedes $h\in G$ ist $h\cdot y=h\cdot(g\cdot x)=(hg)\cdot x\in\operatorname{Orb}_G(x)$, also $\operatorname{Orb}_G(y)\subseteq\operatorname{Orb}_G(x)$. Umgekehrt: $x=g^{-1}\cdot y$, also $\operatorname{Orb}_G(x)\subseteq\operatorname{Orb}_G(y)$. Damit $\operatorname{Orb}_G(x)=\operatorname{Orb}_G(y)$.
“Nur-wenn”: Trivial, denn $y\in\operatorname{Orb}_G(y)=\operatorname{Orb}_G(x)$.

\smallskip
\textbf{(3) Bahnen–Nebenklassen-Bijektion.}
\emph{Wohldefiniertheit:} Ist $[g]=[g']$ in $G/\operatorname{Stab}_G(x)$, so $g'=g s$ mit $s\in\operatorname{Stab}_G(x)$. Dann $g'\cdot x=(g s)\cdot x=g\cdot(s\cdot x)=g\cdot x$.
\emph{Surjektivität:} Jedes $y\in\operatorname{Orb}_G(x)$ hat die Form $y=g\cdot x$; dann ist $\phi([g])=y$.
\emph{Injektivität:} $\phi([g_1])=\phi([g_2])\Rightarrow g_1\cdot x=g_2\cdot x\Rightarrow g_2^{-1}g_1\cdot x=x$, also $g_2^{-1}g_1\in\operatorname{Stab}_G(x)$, d.\,h. $[g_1]=[g_2]$.

\smallskip
\textbf{(4) Konjugation der Stabilisatoren entlang der Bahn.}
Sei $y=t\cdot x$. Dann:
\[
g\in\operatorname{Stab}_G(y)
\ \Longleftrightarrow\ 
g\cdot (t\cdot x)=t\cdot x
\ \Longleftrightarrow\ 
t^{-1} g t \cdot x = x
\ \Longleftrightarrow\
t^{-1} g t \in \operatorname{Stab}_G(x).
\]
Also $g\in t\,\operatorname{Stab}_G(x)\,t^{-1}$ und damit $\operatorname{Stab}_G(y)=t\,\operatorname{Stab}_G(x)\,t^{-1}$.
\end{PROOF}




Definition 13. Sei $G$ eine topologische Gruppe, $X$ ein topologischer Raum. Eine Operation $\mu: G \times X \rightarrow X$ heisst stetige Gruppenoperation, wenn $\mu$ stetig ist.



\begin{DEF}{LIE-B26-07-04}{Stetige Gruppenoperation}
Sei $G$ eine \textbf{topologische Gruppe} und $X$ ein \textbf{topologischer Raum}.  
Eine Abbildung
\[
\mu : G \times X \longrightarrow X, \qquad (g, x) \longmapsto g \cdot x,
\]
heißt \textbf{stetige Gruppenoperation}, wenn gilt:
\begin{enumerate}
    \item $\mu$ ist eine Gruppenoperation, d.\,h.
    \[
    (g h)\cdot x = g\cdot(h\cdot x), \qquad e\cdot x = x \quad \text{für alle } g,h \in G,\, x \in X,
    \]
    wobei $e$ das neutrale Element von $G$ ist;
    \item $\mu$ ist \textbf{stetig} bezüglich der Produkttopologie auf $G\times X$ und der gegebenen Topologie auf $X$.
\end{enumerate}
In diesem Fall sagt man: \emph{$G$ operiert stetig auf $X$} oder \emph{$(G,X)$ ist ein topologisches Wirkungs­paar}.
\end{DEF}



\begin{EXA}{LIE-B26-07-05}{Darstellungen von topologischen Gruppen}
Sei $G$ eine topologische Gruppe und 
\[
\phi : G \longrightarrow \mathrm{GL}(n,\mathbb{K})
\]
ein \textbf{stetiger Gruppenhomomorphismus}.  
Dann induziert $\phi$ eine stetige Gruppenoperation
\[
G \times \mathbb{K}^{n} \longrightarrow \mathbb{K}^{n}, 
\qquad (g, x) \longmapsto \phi(g)\,x.
\]
Diese Operation heißt eine \textbf{Darstellung} von $G$ auf dem Vektorraum 
$V := \mathbb{K}^{n}$.

\smallskip
Man nennt $(V,\phi)$ auch ein \textbf{Darstellungspaar} von $G$.  
Die Abbildung $\phi$ beschreibt, wie die Gruppenstruktur von $G$ 
durch lineare Abbildungen auf $V$ „realisiert“ wird.
\end{EXA}



Beispiel: Sei $\phi: G \rightarrow \mathrm{GL}(n, \mathbb{K})$ ein stetiger Gruppenhomomorphismus. Die induzierte Operation $(g, x) \mapsto \phi(g) x, g \in G, x \in \mathbb{K}^{n}$ heisst Darstellung von $G$ auf $V$. Die Darstellungstheorie von Liegruppen ist ein sehr tiefes und spannendes Forschungsgebiet.




\pagebreak



\section{Die Exponentialabbildung}
Zur Erinnerung: die Taylorreihe der Exponentialfunktion ist

$$
\operatorname{Exp} x=\sum_{k=0}^{\infty} \frac{x^{k}}{k!},
$$

ihr Konvergenzradius ist $+\infty$. Die Taylorreihe der Funktion $\log (x):= \log (1+x)$ ist

$$
\log (x)=\sum_{i=1}^{\infty} \frac{(-1)^{i+1}}{i} x^{i},
$$

ihr Konvergenzradius ist 1 . Insbesondere bekommt man für $x \approx 1$

$$
\log (x)=\log (1+(x-1))=\sum_{i=1}^{\infty} \frac{(-1)^{i+1}}{i}(x-1)^{i} .
$$

Sei wieder $\mathbb{K}=\mathbb{R}$ oder $\mathbb{C}$.




\begin{DEF}{LIE-B26-02-01}{Matrix-Exponential und -Logarithmus}
Sei $\mathbb{K} \in \{\mathbb{R}, \mathbb{C}\}$ und $M_n(\mathbb{K})$ der Raum der $n \times n$-Matrizen.
\begin{enumerate}
\item \textbf{Matrix-Exponential.} Für $A \in M_n(\mathbb{K})$ definieren wir
\[
\exp(A) := \sum_{k=0}^{\infty} \frac{A^k}{k!}.
\]
Diese Reihe konvergiert absolut (bzgl. jeder submultiplikativen Norm) und liefert eine glatte Abbildung
\[
\exp : M_n(\mathbb{K}) \to \mathrm{GL}_n(\mathbb{K}).
\]

\item \textbf{Matrix-Logarithmus (lokal um die Eins).}  
Für $B \in M_n(\mathbb{K})$ mit $\|B\| < 1$ definieren wir
\[
\log(I + B) := \sum_{i=1}^{\infty} \frac{(-1)^{i+1}}{i} B^i.
\]
Die Reihe konvergiert absolut. Damit ist der \emph{Matrix-Logarithmus}
\[
\log : \{ X \in \mathrm{GL}_n(\mathbb{K}) : \|X - I\| < 1 \} \to M_n(\mathbb{K}),
\quad X \mapsto \log(X)
\]
eine glatte Abbildung, gegeben durch die obige Potenzreihe.
\end{enumerate}
\end{DEF}



\begin{PROP}{LIE-B26-02-02}{Grundlegende Eigenschaften von $\exp$ und $\log$}
Für $A, B \in M_n(\mathbb{K})$ gelten:
\begin{enumerate}
\item \textbf{Analytizität:} Die Reihe für $\exp$ konvergiert für alle $A$, also ist $\exp$ auf $M_n(\mathbb{K})$ wohldefiniert, analytisch und $\exp(0)=I$.
\item \textbf{Kommutierender Fall:} Falls $AB = BA$, dann gilt
\[
\exp(A + B) = \exp(A)\exp(B) = \exp(B)\exp(A).
\]
\item \textbf{Ähnlichkeitsinvarianz:} Für $S \in \mathrm{GL}_n(\mathbb{K})$ gilt
\[
\exp(SAS^{-1}) = S\,\exp(A)\,S^{-1}.
\]
\item \textbf{Determinante und Spur:}
\[
\det(\exp A) = e^{\operatorname{tr}(A)}.
\]
\item \textbf{Differentialgleichung:} Für $t \in \mathbb{R}$ ist $\alpha(t) := \exp(tA)$ die eindeutige Lösung der linearen DGL
\[
\alpha'(t) = A\,\alpha(t), \qquad \alpha(0) = I.
\]
\item \textbf{Lokale Inversen:}  
Es existieren offene Umgebungen $U \subset M_n(\mathbb{K})$ von $0$ und $V \subset \mathrm{GL}_n(\mathbb{K})$ von $I$ mit
\[
\exp : U \to V
\]
ein Diffeomorphismus und $\log : V \to U$ die glatte Inverse.
\item \textbf{Ableitung von $\exp$:} Für $A,H \in M_n(\mathbb{K})$ gilt
\[
D(\exp)_A[H] = \int_0^1 \exp((1-s)A)\, H\, \exp(sA)\, ds.
\]
Insbesondere $D(\exp)_0[H] = H$.
\item \textbf{Kommutatorformel:} Für $t \in \mathbb{R}$ gilt
\[
\exp(tA)\,B\,\exp(-tA) = \sum_{k=0}^{\infty} \frac{t^k}{k!}\,(\operatorname{ad}_A)^k(B),
\qquad \operatorname{ad}_A(B) := [A,B] = AB - BA.
\]
\end{enumerate}
\end{PROP}



\begin{PROP}{LIE-B26-02-03}{Matrix-Logarithmus: Eindeutigkeit und Spektralbedingung}
\begin{enumerate}
\item \textbf{Eindeutigkeit nahe $I$:} Gibt es eine Norm mit $\|X - I\| < 1$, so ist $\log(X)$ eindeutig bestimmt und die einzige Lösung von $\exp(A) = X$ mit kleinem $\|A\|$.
\item \textbf{Spektralkriterium:}  
Ist $X \in \mathrm{GL}_n(\mathbb{C})$ ohne Eigenwerte auf $(-\infty,0]$, so existiert ein (holomorph gewählter) \emph{Hauptlogarithmus} $\mathrm{Log}(X)$ mit $\exp(\mathrm{Log}(X)) = X$ und $\sigma(\mathrm{Log}(X)) \subset \{ z \in \mathbb{C} : -\pi < \Im z < \pi \}$.
\item \textbf{Reeller Fall:} Für $X \in \mathrm{GL}_n(\mathbb{R})$ mit $\|X - I\| < 1$ ist $\log(X)$ reell und durch die reelle Potenzreihe definiert.
\end{enumerate}
\end{PROP}


{Bezug zur Lie-Gruppen-Theorie}
\begin{enumerate}
\item Für eine lineare Lie-Gruppe $G \subset \mathrm{GL}_n(\mathbb{K})$ gilt $\exp(\mathfrak{g}) \subset G$, wobei $\mathfrak{g} = T_I G$ die Lie-Algebra ist. Die Kurve $t \mapsto \exp(tA)$ ($A \in \mathfrak{g}$) ist eine \emph{eindimensionale Untergruppe} von $G$.
\item Für $G = \mathrm{GL}_n(\mathbb{K})$ stimmen die Matrixreihen exakt mit der Exponential- bzw. Logarithmusabbildung der Lie-Gruppe überein; insbesondere gilt $\det(\exp A) = e^{\operatorname{tr}(A)}$.
\end{enumerate}






Definition 14. (1) Sei $A \in M_{n}(\mathbb{K})$. Setze

$$
\exp (A)=\sum_{k=0}^{\infty} \frac{1}{k!} A^{k} .
$$

(2) Sei $A \in M_{n}(\mathbb{K})$ mit $\|A-I\|<1$. Setze

$$
\log (A):=\sum_{i=1}^{\infty} \frac{(-1)^{i+1}}{i} A^{i} .
$$





\begin{LEM}{MAT-L03-05}{Grundlegende Rechenregeln der Matrix-Exponentialfunktion}
Sei $A \in M_n(\mathbb{K})$ und $u,v \in \mathbb{K}$. Dann gilt:
\begin{enumerate}
\item $\exp\big((u+v)A\big) = \exp(uA)\,\exp(vA)$.
\item $\exp(A)^{-1} = \exp(-A)$.
\item Seien $A,B\in M_n(\mathbb{K})$ mit $AB=BA$. Dann gilt
\[
\exp(A+B)\;=\;\exp(A)\,\exp(B).
\]
\item Ist $A\in \mathrm{GL}_n(\mathbb{K})$ mit $\|A-I\|<1$, so gilt
\[
\exp(\log A)=A,
\]
wobei $\log A:=\sum_{i=1}^{\infty}\frac{(-1)^{i+1}}{i}(A-I)^i$.

Ist $A\in M_n(\mathbb{K})$ so, dass $\|\exp(A)-I\|<1$, dann gilt
\[
\log(\exp A)=A,
\]
wobei $\log$ durch die Reihe in (1) auf der Umgebung $\{X\in \mathrm{GL}_n(\mathbb{K}) : \|X-I\|<1\}$ definiert ist.
\end{enumerate}

\textit{Beweis.}
\begin{enumerate}
\item
Da die Reihen für $\exp(uA)$ und $\exp(vA)$ absolut konvergieren, darf man die Glieder beliebig umordnen:
\[
\begin{aligned}
\exp(uA)\,\exp(vA)
&= \left(\sum_{i=0}^{\infty} \frac{u^i}{i!} A^i \right)
   \left(\sum_{j=0}^{\infty} \frac{v^j}{j!} A^j \right)
= \sum_{i,j=0}^{\infty} \frac{u^i v^j}{i!\,j!} A^{i+j} \\
&= \sum_{k=0}^{\infty} \left( \sum_{r=0}^k \frac{u^r v^{k-r}}{r!\,(k-r)!} \right) A^k
= \sum_{k=0}^{\infty} \frac{(u+v)^k}{k!} A^k
= \exp((u+v)A).
\end{aligned}
\]
\item
Die Aussage folgt aus (1) durch Wahl von $v=-u$:
\[
\exp(uA)\exp(-uA)=\exp(0)=I
\quad \Rightarrow \quad
\exp(A)^{-1} = \exp(-A).
\]

\textit{Folgerung.}
Damit ist die Abbildung
\[
\exp : M_n(\mathbb{K}) \longrightarrow \mathrm{GL}_n(\mathbb{K})
\]
wohldefiniert und bildet den Raum aller Matrizen in die allgemeine lineare Gruppe ab.

Zur Definition des Matrixlogarithmus schreiben wir für $A \in M_n(\mathbb{K})$ mit $\|A - I\| < 1$:
\[
\log A := \log(A - I) = \sum_{k=1}^{\infty} \frac{(-1)^{k+1}}{k} (A - I)^k.
\]

\item Wegen absoluter Konvergenz der Reihen darf umgruppiert werden:
\[
\begin{aligned}
\exp(A)\exp(B)
&=\left(\sum_{i=0}^{\infty}\frac{A^i}{i!}\right)
  \left(\sum_{j=0}^{\infty}\frac{B^j}{j!}\right)
= \sum_{i=0}^{\infty}\sum_{j=0}^{\infty}\frac{A^i B^j}{i!\,j!}.
\end{aligned}
\]
Da $A$ und $B$ kommutieren, gilt für jedes $k\in\mathbb{N}_0$ der binomiale Lehrsatz
\[
(A+B)^k=\sum_{r=0}^{k}\binom{k}{r}A^{\,r}B^{\,k-r}.
\]
Somit
\[
\sum_{i,j=0}^{\infty}\frac{A^i B^j}{i!\,j!}
=\sum_{k=0}^{\infty}\sum_{r=0}^{k}\frac{1}{r!\,(k-r)!}\,A^{\,r}B^{\,k-r}
=\sum_{k=0}^{\infty}\frac{1}{k!}\sum_{r=0}^{k}\binom{k}{r}A^{\,r}B^{\,k-r}
=\sum_{k=0}^{\infty}\frac{(A+B)^k}{k!}.
\]
Die rechte Seite ist per Definition $\exp(A+B)$.


\item \textbf{Schritt 0 (Vorbereitung).}  
Seien $f(z)=\log(1+z)=\sum_{i=1}^{\infty}\frac{(-1)^{i+1}}{i}z^i$ (Konvergenz für $|z|<1$)
und $g(z)=\exp z=\sum_{m=0}^{\infty}\frac{z^m}{m!}$ (Konvergenz für alle $z$).  
Dann gelten für reelle/komplexe Skalare die Identitäten
\[
g(f(z))=1+z \quad (|z|<1),
\qquad
f(g(z)-1)=z \quad (|g(z)-1|<1).
\]
Da diese Gleichungen auf einer Umgebung des Nullpunkts als Funktionen übereinstimmen,
stimmen ihre Potenzreihen \emph{als formale Reihen} koeffizientenweise überein.

\smallskip
\textbf{Zu (1).}  
Setze $B:=A-I$. Wegen $\|B\|<1$ konvergiert die Reihe
$$\log(I+B)=\sum_{i\ge1}\frac{(-1)^{i+1}}{i}B^i$$ absolut.  
Damit konvergiert auch die zusammengesetzte Matrixreihe
\[
\exp(\log(I+B))=\sum_{m=0}^{\infty}\frac{1}{m!}\Big(\sum_{i=1}^{\infty}\frac{(-1)^{i+1}}{i}B^i\Big)^{\! m},
\]
und absolute Konvergenz erlaubt Umordnen/Assoziieren der Glieder.  
Da $g\circ f=1+z$ als formale Reihe gilt, erhält man nach Ersetzen von $z$ durch $B$
(und Nutzung der gleichen Cauchy-Produktregeln im Matrizenfall) die Identität
\[
\exp(\log(I+B))=I+B=A.
\]

\smallskip
\textbf{Zu (2).}  
Angenommen $\|\exp(A)-I\|<1$. Dann ist $X:=\exp(A)$ in der Definitionsmenge der Logarithmusreihe,
und wir können schreiben
\[
\log(\exp A)=\sum_{i=1}^{\infty}\frac{(-1)^{i+1}}{i}\,(\exp A - I)^i.
\]
Auch hier konvergiert die Reihe absolut.  
Da $f(g(z)-1)=z$ als formale Reihe gilt, folgt durch Ersetzen von $z$ durch die Matrix $A$ (und
absolutes Konvergieren der beteiligten Reihen), dass
\[
\log(\exp A)=A.
\]
\end{enumerate}

\end{LEM}




Lemma 15. Für $A \in M_{n}(\mathbb{K}), u, v \in \mathbb{K}$ gilt\\
(1) $\exp ((u+v) A)=\exp (u A) \exp (v A)$.\\
(2) $\exp (A)^{-1}=\exp (-A)$.

Beweis. (1) Die folgenden Reihen sind absolut konvergent, also darf beliebig umgeordnet werden. Es ist

$$
\begin{aligned}
\operatorname{Exp}(u A) \operatorname{Exp}(v A) & =\left(\sum_{i=0}^{\infty} \frac{u^{i}}{i!} A^{i}\right) \cdot\left(\sum_{j=0}^{\infty} \frac{v^{j}}{i!} A^{j}\right) \\
& =\sum_{i, j=0}^{\infty} \frac{u^{i} v^{j}}{i!j!} A^{i+j} \\
& =\sum_{k=0}^{\infty} \sum_{r=0}^{k} \frac{u^{r} v^{k-r}}{r!(k-r)!} A^{k} \\
& =\sum_{k=0}^{\infty} \frac{(u+v)^{k}}{k!} A^{k} \\
& =\operatorname{Exp}((u+v) A)
\end{aligned}
$$

(2) Folgt aus i).

Insbesondere gilt

$$
\exp : M_{n}(\mathbb{K}) \rightarrow \mathrm{GL}_{n}(\mathbb{K})
$$

Wir setzen

$$
\log A:=\log (A-I)
$$

Lemma 16. Ist $A B=B A$, so gilt $\exp (A+B)=\exp (A) \exp (B)$.\\
Beweis.

$$
\begin{aligned}
\exp (A) \exp (B) & =\sum_{i=0}^{\infty} \frac{1}{i!} A^{i} \sum_{j=0}^{\infty} \frac{1}{j!} B^{j} \\
& =\sum_{k=0}^{\infty} \sum_{r=0}^{k} \frac{1}{k!}\binom{k}{r} A^{r} B^{k-r} \\
& =\sum_{k=0}^{\infty} \frac{1}{k!}(A+B)^{k} \\
& =\exp (A+B)
\end{aligned}
$$

Lemma 17. (1) Ist $\|A-I\|<1$, so ist $\exp \log A=A$.\\
(2) Ist $\|\exp A-I\|<1$, so ist $\log \exp A=A$.

Beweis. Die folgenden Gleichungen sind Gleichungen für formale Potenzreihen:

$$
\begin{aligned}
\sum_{m=0}^{\infty} \frac{1}{m!}\left(\sum_{i=1}^{\infty} \frac{(-1)^{i+1}}{i}(X-1)^{i}\right)^{m} & =X \\
\sum_{i=1}^{\infty} \frac{(-1)^{i+1}}{i}\left(\sum_{m=0}^{\infty} \frac{1}{m} X^{m}\right)^{i} & =X
\end{aligned}
$$

Für reelle $X \approx 1$ gelten diese Gleichungen, also stimmen alle Koeffizienten überein und diese Gleichungen gelten auch für beliebige $X$. Ersetzt man $X$ durch $A$, so folgen die Behauptungen.



Korollar 18. 

Die Einschränkung von $\exp$ auf
$$
N_{\left.M_{n}(\mathbb{K})\right)}(0, \log 2)=\left\{A \in M_{n}(\mathbb{K}):\|A\|<\log 2\right\}
$$
ist ein lokaler Diffeomorphismus mit lokaler Inverser $\log$.


Beweis. exp und log sind unendlich oft differenzierbar, da Potenzreihen. Es gilt

$$
\|\exp (A)-I\| \leq \sum_{i=1}^{\infty} \frac{1}{i!}\left\|A^{i}\right\| \leq \sum_{i=1}^{\infty} \frac{1}{i!}\|A\|^{i}=e^{\|A\|}-1
$$

also ist $\|\exp (A)-I\|<1$ falls $\|A\|<\log 2$.\\



\begin{KORO}{LIE-B26-02-04}{Lokaler Diffeomorphismus der Exponentialabbildung}
Sei $\|\cdot\|$ eine submultiplikative Matrixnorm auf $M_n(\mathbb{K})$ und
\[
N := N_{M_n(\mathbb{K})}(0,\log 2)=\{A\in M_n(\mathbb{K}): \|A\|<\log 2\}.
\]
Dann ist die Einschränkung
\[
\exp\big|_{N} : N \longrightarrow \mathrm{GL}_n(\mathbb{K})
\]
ein (auf sein Bild) lokaler Diffeomorphismus mit lokaler Inversen $\log$; genauer gilt
\[
\exp(N)\subset U:=\{X\in \mathrm{GL}_n(\mathbb{K}) : \|X-I\|<1\},
\quad\text{und}\quad
\log:U\to M_n(\mathbb{K})
\]
ist glatt sowie $\log\circ\exp=\mathrm{id}_{N}$ und $\exp\circ\log=\mathrm{id}_{U}$.
\end{KORO}


\begin{PROOF}{LIE-B26-02-05}{P: Lokaler Diffeomorphismus der Exponentialabbildung}
(1) \textit{Bild liegt in $U$.} Für $A\in M_n(\mathbb{K})$ gilt wegen submultiplikativität
\[
\|\exp(A)-I\| \le \sum_{k=1}^\infty \frac{\|A\|^k}{k!} = e^{\|A\|}-1.
\]
Insbesondere ist für $\|A\|<\log 2$ die rechte Seite $<1$, also $\exp(A)\in U$ und damit $\exp(N)\subset U$.

(2) \textit{Glatte Inversenbeziehung.} Nach Lemma~\textit{(Wechselseitigkeit von $\exp$ und $\log$ im Konvergenzbereich)} gilt
\[
\log(\exp A)=A \quad\text{für }\|A\|<\log 2,
\qquad
\exp(\log X)=X \quad\text{für }\|X-I\|<1.
\]
Damit sind $\log$ und $\exp$ auf $U$ bzw. $N$ wechselseitige Inversen. Beide sind als Potenzreihen glatt.

(3) \textit{Lokaler Diffeomorphismus.} Aus (2) folgt bereits, dass $\exp\big|_{N}$ ein Diffeomorphismus zwischen $N$ und $U$ ist (insbesondere ein lokaler Diffeomorphismus). Alternativ: $T_{0}\exp=\mathrm{id}$ (Theorem über die Ableitung der Exponentialabbildung), so dass der lokale Diffeomorphismus am Ursprung auch unmittelbar aus dem inversen Funktionensatz folgt; Schritt (2) liefert explizit die Inverse $\log$ auf $U$.
\end{PROOF}





Satz 19. Für $A, B \in M_{n}(\mathbb{C})$ gilt

$$
\left.\frac{d}{d t}\right|_{t=0} \exp (A+t B)=F(\operatorname{ad} A)(B) \exp (A)
$$

wobei

$$
\operatorname{ad} A: M_{n}(\mathbb{C}) \rightarrow M_{n}(\mathbb{C}), C \mapsto A C-C A=[A, C]
$$

die adjungierte Operation ist und

$$
F(t)=\sum_{r=0}^{\infty} \frac{1}{(r+1)!} t^{r}=\frac{e^{t}-1}{t}
$$

Speziell gilt

$$
\left.\frac{d}{d t}\right|_{t=0} \exp (t B)=B
$$

Beweis. 

\begin{PROOF}{LIE-B26-02-08}{P:(2) Ableitung der Matrix-Exponentialfunktion entlang einer Richtung}
Per Induktion zeigt man leicht
$$
\sum_{k=0}^{m-1}\binom{k}{i}=\binom{m}{i+1}
$$
Per Induktion zeit man weiterhin
$$
A^{k} B=\sum_{i=0}^{k}\binom{k}{i}(\operatorname{ad} A)^{i} B \cdot A^{k-i}
$$
Der Induktionsanfang $k=0$ ist trivial. Induktionsschritt:
$$
\begin{aligned}
A^{k+1} B & =A\left(A^{k} B\right) \\
& =A \sum_{i=0}^{k}\binom{k}{i}(\operatorname{ad} A)^{i} B \cdot A^{k-i} \\
& =\sum_{i=0}^{k}\binom{k}{i}\left[\left(A \cdot(\operatorname{ad} A)^{i} B-(\operatorname{ad} A)^{i} B A\right) \cdot A^{k-i}+(\operatorname{ad} A)^{i} B \cdot A^{k+1-i}\right] \\
& =\sum_{i=0}^{k}\binom{k}{i}\left[(\operatorname{ad} A)^{i+1} B A^{k-i}+(\operatorname{ad} A)^{i} B A^{k+1-i}\right] \\
& =\sum_{i=1}^{k+1}\left[\binom{k}{i-1}(\operatorname{ad} A)^{i} B A^{k+1-i}+\binom{k}{i}(\operatorname{ad} A)^{i} B A^{k+1-i}\right]+B A^{k+1} \\
& =\sum_{i=1}^{k+1}\binom{k+1}{i}(\operatorname{ad} A)^{i} B A^{k+1-i}+B A^{k+1} \\
& =\sum_{i=0}^{k+1}\binom{k+1}{i}(\operatorname{ad} A)^{i} B A^{k+1-i} .
\end{aligned}
$$
Wegen
$$
\left.\frac{d}{d t}\right|_{t=0}(A+t B)^{m}=\sum_{k=0}^{m-1} A^{k} B A^{m-1-k}
$$
folgt
$$
\begin{aligned}
\left.\frac{d}{d t}\right|_{t=0} \exp (A+t B) & =\left.\sum_{m=0}^{\infty} \frac{1}{m!} \frac{d}{d t}\right|_{t=0}(A+t B)^{m} \\
& =\sum_{m=0}^{\infty} \frac{1}{m!} \sum_{k=0}^{m-1} \sum_{i=0}^{k}\binom{k}{i}(\operatorname{ad} A)^{i} B A^{m-i-1} \\
& =\sum_{m=0}^{\infty} \frac{1}{m!} \sum_{i=0}^{m-1}\left(\sum_{k=i}^{m-1}\binom{k}{i}\right)(\operatorname{ad} A)^{i} B A^{m-i-1} \\
& =\sum_{m=0}^{\infty} \frac{1}{m!} \sum_{i=0}^{m-1}\binom{m}{i+1}(\operatorname{ad} A)^{i} B A^{m-1-i} \\
& =\sum_{i=0}^{\infty} \frac{(\operatorname{ad} A)^{i} B}{(i+1)!} \sum_{m=i+1}^{\infty} \frac{1}{(m-i-1)!} A^{m-i-1} \\
& =F(\operatorname{ad} A) B \cdot \exp A
\end{aligned}
$$
\end{PROOF}



\begin{THEO}{LIE-B26-02-06}{Ableitung der Matrix-Exponentialfunktion entlang einer Richtung}
Seien $A,B \in M_{n}(\mathbb{C})$. Dann gilt
\[
\left.\frac{d}{dt}\right|_{t=0} \exp(A+tB)
 \;=\; F(\operatorname{ad}A)(B)\,\exp(A),
\]
wobei $\operatorname{ad}A : M_n(\mathbb{C}) \to M_n(\mathbb{C}),\; C \mapsto [A,C]=AC-CA$ 
die adjungierte Operation ist und
\[
F(t) \;=\; \sum_{r=0}^{\infty} \frac{t^{r}}{(r+1)!} \;=\; \frac{e^{t}-1}{t}.
\]
Insbesondere gilt
\[
\left.\frac{d}{dt}\right|_{t=0} \exp(tB) \;=\; B.
\]
\end{THEO}



\begin{PROOF}{LIE-B26-02-07}{P: Ableitung der Matrix-Exponentialfunktion entlang einer Richtung}
Zunächst zeigen wir zwei elementare Summenidentitäten per Induktion:
\begin{enumerate}\itemsep0.2em
\item[(S1)] Für $m\in\mathbb{N}$ und $i\in\{0,\dots,m-1\}$ gilt 
\[
\sum_{k=0}^{m-1} \binom{k}{i} \;=\; \binom{m}{i+1}.
\]
\item[(S2)] Für $k\in\mathbb{N}_0$ gilt
\[
A^{k}B \;=\; \sum_{i=0}^{k} \binom{k}{i} (\operatorname{ad}A)^{i}(B)\; A^{k-i}.
\]
\end{enumerate}
Der Induktionsanfang von (S2) ist klar für $k=0$. Für den Schritt $k\mapsto k+1$ verwenden wir
\[
\begin{aligned}
A^{k+1}B &= A(A^{k}B)
= A \sum_{i=0}^{k} \binom{k}{i} (\operatorname{ad}A)^{i}(B)\,A^{k-i}\\
&=\sum_{i=0}^{k} \binom{k}{i} 
\Big( \underbrace{A(\operatorname{ad}A)^{i}(B) - (\operatorname{ad}A)^{i}(B)A}_{(\operatorname{ad}A)^{i+1}(B)} \Big) A^{k-i}
\;+\; \sum_{i=0}^{k} \binom{k}{i} (\operatorname{ad}A)^{i}(B)\,A^{k+1-i}\\
&=\sum_{i=1}^{k+1} \bigg[ \binom{k}{i-1} + \binom{k}{i} \bigg] (\operatorname{ad}A)^{i}(B) A^{k+1-i} \;+\; (\operatorname{ad}A)^{0}(B)A^{k+1}\\
&=\sum_{i=0}^{k+1} \binom{k+1}{i} (\operatorname{ad}A)^{i}(B) A^{k+1-i},
\end{aligned}
\]
womit (S2) gezeigt ist. Die Identität (S1) ist der bekannte Summationssatz für Binomialkoeffizienten.

Nun leiten wir die Exponentialreihe gliedweise ab (absolut konvergent in jeder Matrixnorm):
\[
\left.\frac{d}{dt}\right|_{t=0} (A+tB)^{m}
= \sum_{k=0}^{m-1} A^{k} B A^{m-1-k}.
\]
Mit (S2) erhalten wir für jedes $k$
\[
A^{k} B = \sum_{i=0}^{k} \binom{k}{i} (\operatorname{ad}A)^{i}(B) A^{k-i}.
\]
Somit
\[
\begin{aligned}
\left.\frac{d}{dt}\right|_{t=0} \exp(A+tB)
&= \sum_{m=0}^{\infty} \frac{1}{m!} \left.\frac{d}{dt}\right|_{t=0} (A+tB)^{m}
= \sum_{m=0}^{\infty} \frac{1}{m!} \sum_{k=0}^{m-1} A^{k} B A^{m-1-k}\\
&= \sum_{m=0}^{\infty} \frac{1}{m!} \sum_{k=0}^{m-1} \sum_{i=0}^{k} \binom{k}{i} (\operatorname{ad}A)^{i}(B) A^{m-1-i}\\
&= \sum_{m=0}^{\infty} \frac{1}{m!} \sum_{i=0}^{m-1} \Big(\sum_{k=i}^{m-1} \binom{k}{i}\Big) (\operatorname{ad}A)^{i}(B) A^{m-1-i}\\
&\stackrel{(S1)}{=} \sum_{m=0}^{\infty} \frac{1}{m!} \sum_{i=0}^{m-1} \binom{m}{i+1} (\operatorname{ad}A)^{i}(B) A^{m-1-i}.
\end{aligned}
\]
Vertauschen der Summation (absolut konvergent) liefert
\[
\begin{aligned}
\left.\frac{d}{dt}\right|_{t=0} \exp(A+tB)
&= \sum_{i=0}^{\infty} \frac{(\operatorname{ad}A)^{i}(B)}{(i+1)!} \sum_{m=i+1}^{\infty} \frac{A^{m-1-i}}{(m-1-i)!}\\
&= \Big(\sum_{i=0}^{\infty} \frac{(\operatorname{ad}A)^{i}}{(i+1)!}\Big)(B) \cdot \Big(\sum_{r=0}^{\infty} \frac{A^{r}}{r!}\Big)\\
&= F(\operatorname{ad}A)(B)\,\exp(A),
\end{aligned}
\]
wobei $F(t)=\sum_{r=0}^{\infty}\frac{t^{r}}{(r+1)!}=\frac{e^{t}-1}{t}$.  
Für den Spezialfall $A=0$ folgt sofort
\[
\left.\frac{d}{dt}\right|_{t=0} \exp(tB) = B.
\]
\end{PROOF}


Satz 20. Seien $A, C \in M_{n}(\mathbb{R}), a<0<b$. Dann hat die Differentialgleichung

$$
\begin{aligned}
\alpha^{\prime}(t) & =\alpha(t) A \\
\alpha(0) & =C
\end{aligned}
$$

genau eine Lösung $\alpha:(a, b) \rightarrow M_{n}(\mathbb{R})$, nämlich $\alpha(t)=C \exp (t A)$.


Beweis. Es ist

$$
\begin{aligned}
\alpha^{\prime}(t) & =C \sum_{m=0}^{\infty} \frac{d}{d t} \frac{(t A)^{m}}{m!} \\
& =C \sum_{m=1}^{\infty} \frac{t^{m-1} A^{m}}{(m-1)!} \\
& =C \sum_{m=0}^{\infty} \frac{t^{m} A^{m-1}}{m!} A \\
& =\alpha(t) A
\end{aligned}
$$

Ist $\beta(t)$ eine weitere Lösung, so gilt mit $\gamma(t)=\beta(t) \exp (-t A)$

$$
\begin{aligned}
\gamma^{\prime}(t) & =\beta^{\prime}(t) \exp (-t A)-\beta(t) \exp (-t A) A \\
& =\beta(t) A \exp (-t A)-\beta(t) \exp (-t A) A \\
& =0
\end{aligned}
$$

also ist $\gamma$ konstant. Wegen $\gamma(0)=\beta(0)=C$ folgt $\beta(t)=C \exp (t A)= \alpha(t)$.



\begin{THEO}{LIE-B26-02-09}{Lineare Matrix-DGL mit konstanter Koeffizientenmatrix}
Seien $A,C\in M_n(\mathbb{R})$ und $a<0<b$. Die Anfangswertaufgabe
\[
\alpha'(t)=\alpha(t)\,A,\qquad \alpha(0)=C,
\]
besitzt auf $(a,b)$ genau eine Lösung, nämlich
\[
\alpha(t)=C\,\exp(tA).
\]
\end{THEO}



\begin{PROOF}{LIE-B26-02-10}{P: Lineare Matrix-DGL mit konstanter Koeffizientenmatrix}
\textbf{Existenz.} Definiere $\alpha(t):=C\exp(tA)=C\sum_{m=0}^\infty \frac{(tA)^m}{m!}$. Dann gilt wegen gliedweiser Ableitbarkeit der absolut konvergenten Reihe:
\[
\begin{aligned}
\alpha'(t)
&= C\sum_{m=0}^\infty \frac{d}{dt}\Big(\frac{t^m}{m!}A^m\Big)
= C\sum_{m=1}^\infty \frac{t^{m-1}}{(m-1)!}A^m \\
&= C\sum_{m=0}^\infty \frac{t^{m}}{m!}A^{m}A
= \big(C\sum_{m=0}^\infty \frac{t^{m}}{m!}A^{m}\big)A
= \alpha(t)A,
\end{aligned}
\]
und $\alpha(0)=C\exp(0)=C$.

\medskip
\textbf{Eindeutigkeit.} Sei $\beta:(a,b)\to M_n(\mathbb{R})$ eine weitere Lösung. Setze
\[
\gamma(t):=\beta(t)\exp(-tA).
\]
Dann
\[
\begin{aligned}
\gamma'(t)
&=\beta'(t)\exp(-tA)+\beta(t)\,\frac{d}{dt}\exp(-tA) \\
&=\beta(t)A\exp(-tA)-\beta(t)\exp(-tA)A
=0,
\end{aligned}
\]
da $\frac{d}{dt}\exp(-tA)=-\exp(-tA)A$. Also ist $\gamma$ konstant, mithin
\[
\gamma(t)\equiv \gamma(0)=\beta(0)\exp(0)=C.
\]
Folglich $\beta(t)=\gamma(t)\exp(tA)=C\exp(tA)=\alpha(t)$.
\end{PROOF}




Definition 21. 


Sei $G$ eine Matrixgruppe. Eine 1-Parameter-Semigruppe ist eine stetige Funktion $\gamma:(-\epsilon, \epsilon) \rightarrow G$, die in 0 differenzierbar ist (d.h. $\lim _{t \rightarrow 0} \frac{\gamma(t)-\gamma(0)}{t} \in M_{n}(\mathbb{K})$ existiert) und $\gamma(s+t)=\gamma(s) \gamma(t)$ für alle $s, t, s+t \in(-\epsilon, \epsilon)$ erfüllt. Ist $\epsilon=\infty$, so heisst $\gamma$ eine 1-ParameterUntergruppe von $G$.



\begin{DEF}{LIE-B26-02-11}{1-Parameter-(Semi)gruppe in einer Matrixgruppe}
Sei $G\subset \GL_n(\mathbb{K})$ eine Matrixgruppe und $\epsilon\in(0,\infty]$.
\begin{enumerate}\itemsep0.2em
\item Eine \emph{1-Parameter-Semigruppe in $G$} ist eine stetige Abbildung
\[
\gamma:(-\epsilon,\epsilon)\longrightarrow G,
\]
die in $0$ differenzierbar ist (d.\,h. der Grenzwert
$\displaystyle\lim_{t\to 0}\frac{\gamma(t)-\gamma(0)}{t}\in M_n(\mathbb{K})$ existiert) und die \emph{Semigruppeneigenschaft}
\[
\gamma(s+t)=\gamma(s)\,\gamma(t)
\quad\text{für alle $s,t$ mit $s,t,s+t\in(-\epsilon,\epsilon)$}
\]
erfüllt.
\item Ist zusätzlich $\epsilon=\infty$ (d.\,h. $\gamma:\mathbb{R}\to G$) und gilt die obige Gleichung für alle $s,t\in\mathbb{R}$, so heißt $\gamma$ eine \emph{1-Parameter-Untergruppe von $G$}.
\end{enumerate}
\end{DEF}




SATZ 22. 

Sei $\gamma:(-\epsilon, \epsilon) \rightarrow G$ eine 1-Parameter-Semigruppe. Dann ist $\gamma$ differenzierbar und $\gamma^{\prime}(t)=\gamma^{\prime}(0) \gamma(t)=\gamma(t) \gamma^{\prime}(0)$.

Beweis.

$$
\gamma^{\prime}(t)=\lim _{h \rightarrow 0} \frac{1}{h}(\gamma(t+h)-\gamma(t))=\lim _{h \rightarrow 0} \frac{1}{h}(\gamma(h)-I d) \gamma(t)=\gamma^{\prime}(0) \gamma(t)
$$



\begin{PROP}{LIE-B26-02-12}{Ableitung von 1-Parameter-Semigruppen}
Sei $G\subset \GL_n(\mathbb{K})$ eine Matrixgruppe und
$\gamma:(-\varepsilon,\varepsilon)\to G$ eine 1-Parameter-Semigruppe,
d.\,h. $\gamma$ ist stetig, in $0$ differenzierbar und erfüllt
$\gamma(s+t)=\gamma(s)\gamma(t)$, sofern $s,t,s+t\in(-\varepsilon,\varepsilon)$.
Dann gilt:
\begin{enumerate}\itemsep0.2em
\item $\gamma$ ist in jedem $t\in(-\varepsilon,\varepsilon)$ differenzierbar.
\item Für alle $t$ gilt
\[
\gamma'(t)=\gamma'(0)\,\gamma(t)=\gamma(t)\,\gamma'(0).
\]
\end{enumerate}
\end{PROP}



\begin{PROOF}{LIE-B26-02-13}{P: Ableitung von 1-Parameter-Semigruppen}
Fixiere $t\in(-\varepsilon,\varepsilon)$. Für hinreichend kleines $h$ ist auch $t+h\in(-\varepsilon,\varepsilon)$ und nach der Semigruppeneigenschaft
\[
\gamma(t+h)=\gamma(t)\gamma(h)=\gamma(h)\gamma(t).
\]
Dann
\[
\frac{\gamma(t+h)-\gamma(t)}{h}
=\frac{\gamma(t)\gamma(h)-\gamma(t)}{h}
=\gamma(t)\,\frac{\gamma(h)-I}{h},
\]
und wegen der Existenz des Grenzwerts $\lim_{h\to 0}\frac{\gamma(h)-I}{h}=\gamma'(0)$ folgt
\[
\gamma'(t)=\lim_{h\to 0}\frac{\gamma(t+h)-\gamma(t)}{h}
=\gamma(t)\,\gamma'(0).
\]
Analog, aus $\gamma(t+h)=\gamma(h)\gamma(t)$ erhalten wir
\[
\frac{\gamma(t+h)-\gamma(t)}{h}
=\frac{\gamma(h)-I}{h}\,\gamma(t)\;\xrightarrow{h\to 0}\;\gamma'(0)\,\gamma(t),
\]
also auch $\gamma'(t)=\gamma'(0)\gamma(t)$. Damit ist $\gamma$ überall differenzierbar und beide Darstellungen der Ableitung gelten.
\end{PROOF}




Satz 23. Sei $\gamma:(-\epsilon, \epsilon) \rightarrow G$ eine 1-Parameter-Semigruppe. Dann lässt sich $\gamma$ eindeutig fortsetzen zu einer 1-Parameter-Untergruppe $\tilde{\gamma}$ : $\mathbb{R} \rightarrow G$.

Beweis. Eindeutigkeit: Angenommen $\tilde{\gamma}$ existiert. Dann gilt für alle $m \in \mathbb{N}, t \in \mathbb{R}$

$$
\tilde{\gamma}(t)=\tilde{\gamma}\left(m \cdot \frac{t}{m}\right)=\tilde{\gamma}\left(\frac{t}{m}\right)^{m}
$$

Ist $m$ genügend groß, so dass $\frac{t}{m} \in(-\epsilon, \epsilon)$, so folgt $\tilde{\gamma}(t)=\gamma\left(\frac{t}{m}\right)^{m}$. Das zeigt die Eindeutigkeit.

Nun zur Existenz. Wir definieren $\tilde{\gamma}(t):=\gamma\left(\frac{t}{m}\right)^{m}$, wobei $m \in \mathbb{N}$ so groß ist, dass $\frac{t}{m} \in(-\epsilon, \epsilon)$. Zu zeigen: das ist wohldefiniert, also unabhängig von der Wahl von $m$. Sei $k \in \mathbb{N}$ mit $\frac{t}{k} \in(-\epsilon, \epsilon)$. Dann gilt

$$
\gamma\left(\frac{t}{k}\right)^{k}=\gamma\left(m \frac{t}{m k}\right)^{k}=\gamma\left(\frac{t}{m k}\right)^{m k}
$$

und

$$
\gamma\left(\frac{t}{m}\right)^{m}=\gamma\left(k \frac{t}{m k}\right)^{m}=\gamma\left(\frac{t}{m k}\right)^{m k} .
$$

Also ist $\gamma\left(\frac{t}{k}\right)^{k}=\gamma\left(\frac{t}{m}\right)^{m}$ und $\tilde{\gamma}$ ist wohldefiniert.\\
Für $t, s \in \mathbb{R}$ wähle $m$ so groß, dass $\frac{t}{m}, \frac{s}{m}, \frac{t+s}{m} \in(-\epsilon, \epsilon)$. Dann ist

$$
\tilde{\gamma}(s+t)=\gamma\left(\frac{t+s}{m}\right)^{m}=\gamma\left(\frac{t}{m}\right)^{m} \gamma\left(\frac{s}{m}\right)^{m}=\tilde{\gamma}(t) \tilde{\gamma}(s)
$$



\begin{PROP}{LIE-B26-02-14}{Eindeutige Fortsetzung einer 1-Parameter-Semigruppe}
Sei $G\subset \GL_n(\mathbb{K})$ eine Matrixgruppe und 
$\gamma:(-\varepsilon,\varepsilon)\to G$ eine 1\,-Parameter-Semigruppe.
Dann existiert genau eine 1\,-Parameter-Untergruppe
\[
\tilde{\gamma}:\mathbb{R}\longrightarrow G,\qquad
\text{stetig mit }\ \tilde{\gamma}(s+t)=\tilde{\gamma}(s)\tilde{\gamma}(t)\ \ (\forall s,t\in\mathbb{R}),
\]
die $\gamma$ auf $(-\varepsilon,\varepsilon)$ fortsetzt.
\end{PROP}



\begin{PROOF}{LIE-B26-02-15}{P: Eindeutige Fortsetzung einer 1-Parameter-Semigruppe}
\textbf{Eindeutigkeit.}
Ist $\tilde{\gamma}$ eine Fortsetzung mit der Homomorphieeigenschaft,
so gilt für $m\in\mathbb{N}$ und $t\in\mathbb{R}$:
\[
\tilde{\gamma}(t)=\tilde{\gamma}\!\left(m\cdot\frac{t}{m}\right)
=\big(\tilde{\gamma}(t/m)\big)^{m}.
\]
Für $m$ so groß, dass $t/m\in(-\varepsilon,\varepsilon)$, folgt
$\tilde{\gamma}(t)=\gamma(t/m)^{m}$. Damit ist $\tilde{\gamma}$ durch $\gamma$
eindeutig bestimmt.

\medskip
\textbf{Existenz / Definition.}
Für $t\in\mathbb{R}$ wähle $m\in\mathbb{N}$ mit $t/m\in(-\varepsilon,\varepsilon)$
und setze
\[
\tilde{\gamma}(t):=\gamma(t/m)^{\,m}.
\]
\emph{Wohldefiniertheit:}
Seien $m,k\in\mathbb{N}$ mit $t/m,t/k\in(-\varepsilon,\varepsilon)$. Dann
\[
\gamma(t/k)^{k}
=\gamma\!\left(m\cdot \frac{t}{mk}\right)^{k}
=\big(\gamma(t/(mk))^{\,m}\big)^{k}
=\gamma(t/(mk))^{\,mk}
=\big(\gamma(t/(mk))^{\,k}\big)^{m}
=\gamma(t/m)^{m}.
\]
Also ist $\tilde{\gamma}$ wohldefiniert.

\medskip
\textbf{Homomorphieeigenschaft.}
Für $s,t\in\mathbb{R}$ wähle $m\in\mathbb{N}$ mit
$t/m,\ s/m,\ (s+t)/m \in(-\varepsilon,\varepsilon)$. Dann
\[
\tilde{\gamma}(s+t)=\gamma\!\left(\tfrac{s+t}{m}\right)^{m}
=\big(\gamma(s/m)\,\gamma(t/m)\big)^{m}
=\gamma(s/m)^{m}\,\gamma(t/m)^{m}
=\tilde{\gamma}(s)\,\tilde{\gamma}(t),
\]
wobei die Semigruppeneigenschaft von $\gamma$ auf $(-\varepsilon,\varepsilon)$ benutzt wurde.

\medskip
\textbf{Stetigkeit.}
Stetigkeit an $0$: Sei $t\to 0$. Wähle $m$ so groß, dass 
$t/m\in(-\varepsilon,\varepsilon)$ und $|t|$ so klein, dass ein fixes $m$ genügt.
Dann
\[
\tilde{\gamma}(t)=\gamma(t/m)^{\,m}\to \gamma(0)^{\,m}=I^{m}=I=\tilde{\gamma}(0),
\]
da $\gamma$ in $0$ stetig ist. Für allgemeines $t_0$ nutze Homomorphie:
\[
\tilde{\gamma}(t)=\tilde{\gamma}\big((t-t_0)+t_0\big)
=\tilde{\gamma}(t-t_0)\,\tilde{\gamma}(t_0)\xrightarrow[t\to t_0]{}\tilde{\gamma}(0)\,\tilde{\gamma}(t_0)=\tilde{\gamma}(t_0).
\]
Damit ist $\tilde{\gamma}$ stetig auf $\mathbb{R}$.

\medskip
Somit ist $\tilde{\gamma}:\mathbb{R}\to G$ eine stetige Gruppenhomomorphie
von $(\mathbb{R},+)$ nach $G$ und eine Fortsetzung von $\gamma$, und wegen des ersten Teils eindeutig.
\end{PROOF}



Theorem 24. Jede 1-Parameter-Gruppe $\gamma: \mathbb{R} \rightarrow G$ hat die Form $\gamma(t)=\exp (t A)$ mit $A \in M_{n}(\mathbb{K})$.

Beweis. Ist $A=\gamma^{\prime}(0)$, so erfüllt $\gamma$ die DFL $\gamma^{\prime}(t)=A \gamma(t), \gamma(0)=$ Id. Die eindeutige Lösung ist $\gamma(t)=\exp (t A)$.



\begin{THEO}{LIE-B26-02-16}{Charakterisierung von 1\,-Parameter-Gruppen durch $\exp$-Abbildung}
Sei $G\subset \GL_n(\mathbb{K})$ eine Matrixgruppe. Jede 1\,-Parameter-Gruppe 
$\gamma:\mathbb{R}\to G$ besitzt die Form
\[
\gamma(t)=\exp(tA)\quad\text{für ein eindeutig bestimmtes }A\in M_n(\mathbb{K}).
\]
\end{THEO}



\begin{PROOF}{LIE-B26-02-17}{P: Charakterisierung von 1\,-Parameter-Gruppen durch $\exp$-Abbildung}
Setze $A:=\gamma'(0)\in M_n(\mathbb{K})$, deren Existenz aus der
Gruppenstruktur von $\gamma$ folgt. Nach Proposition \textsc{GRP-LIN-02} gilt für alle $t\in\mathbb{R}$:
\[
\gamma'(t)=\gamma'(0)\,\gamma(t)=A\,\gamma(t),\qquad \gamma(0)=I.
\]
Damit ist $\gamma$ Lösung des linearen Matrix-AWPs
\[
\alpha'(t)=A\,\alpha(t),\qquad \alpha(0)=I.
\]
Die eindeutige (globale) Lösung dieses Anfangswertproblems ist bekanntlich
$\alpha(t)=\exp(tA)$ (vgl.\ Definition der Matrix-Exponentialfunktion und
Satz zur Eindeutigkeit linearer Systeme). Also $\gamma(t)=\exp(tA)$ für alle $t$.

\emph{Eindeutigkeit von $A$:}
Aus $\gamma(t)=\exp(tA)=\exp(tB)$ für alle $t$ folgt durch Ableiten in $t=0$,
dass $A=\gamma'(0)=B$.
\end{PROOF}



\pagebreak
\printbibliography
\end{document}